%File: anonymous-submission-latex-2025.tex
\documentclass[letterpaper]{article} % DO NOT CHANGE THIS
\usepackage{aaai25}  % DO NOT CHANGE THIS
\usepackage{times}  % DO NOT CHANGE THIS
\usepackage{helvet}  % DO NOT CHANGE THIS
\usepackage{courier}  % DO NOT CHANGE THIS
\usepackage[hyphens]{url}  % DO NOT CHANGE THIS
\usepackage{graphicx} % DO NOT CHANGE THIS
\urlstyle{rm} % DO NOT CHANGE THIS
\def\UrlFont{\rm}  % DO NOT CHANGE THIS
\usepackage{natbib}  % DO NOT CHANGE THIS AND DO NOT ADD ANY OPTIONS TO IT
\usepackage{caption} % DO NOT CHANGE THIS AND DO NOT ADD ANY OPTIONS TO IT
\frenchspacing  % DO NOT CHANGE THIS
\setlength{\pdfpagewidth}{8.5in} % DO NOT CHANGE THIS
\setlength{\pdfpageheight}{11in} % DO NOT CHANGE THIS
%
% These are recommended to typeset algorithms but not required. See the subsubsection on algorithms. Remove them if you don't have algorithms in your paper.
\usepackage{algorithm}
% \usepackage{algorithmic}

% \usepackage{hyperref}
\usepackage{url}
\usepackage{booktabs}
% \usepackage{amsmath,amsthm,amsfonts,amssymb,amscd}
\usepackage{amsfonts}
% \usepackage{bbm}
% \usepackage{hyperref}
% \usepackage{lastpage}
\usepackage{enumerate}
% \usepackage{fancyhdr}
\usepackage{mathrsfs}
% \usepackage{graphicx}
\usepackage{listings}
% \usepackage{hyperref}
\usepackage{mathtools}
% \usepackage{algpseudocode}
\usepackage{todonotes}
\usepackage{adjustbox}
% \usepackage{wrapfig}
\usepackage{extarrows}
\usepackage{scrextend}
\usepackage{multirow}
\usepackage{makecell}
\usepackage{pifont}
% \usepackage{microtype}
% \usepackage[subtle]{savetrees}
\usepackage{array}
\usepackage{algpseudocode}
\usepackage{cleveref}

\newcommand{\ba}{\mathbf{a}}
\newcommand{\bc}{\mathbf{c}}
\newcommand{\bx}{\mathbf{x}}
\newcommand{\bz}{\mathbf{z}}
\newcommand{\by}{\mathbf{y}}
\newcommand{\bs}{\mathbf{s}}
\newcommand{\bh}{\mathbf{h}}
\newcommand{\bb}{\mathbf{b}}
\newcommand{\br}{\mathbf{r}}
\newcommand{\be}{\mathbf{e}}
\newcommand\Exp{\mathbb{E}}
\newcommand\KL{\text{KL}}
\newcommand\Lc{\mathcal{L}}
\newcommand\Uc{\mathcal{U}}
\newcommand\Hc{\mathcal{H}}
\newcommand\Ic{\mathcal{I}}
\newcommand\Rc{\mathcal{R}}
\newcommand\MI{\text{MI}}
\newcommand\D{\text{D}}
\newcommand{\btheta}{{\boldsymbol{\theta}}}
\newcommand{\bphi}{{\boldsymbol{\phi}}}
\newcommand{\bepsilon}{{\boldsymbol{\epsilon}}}
\newcommand{\bmu}{{\boldsymbol{\mu}}}
\newcommand{\bnu}{{\boldsymbol{\nu}}}
\newcommand{\bSigma}{{\boldsymbol{\Sigma}}}
\newcommand{\bzero}{\mathbf{0}}
\newcommand{\bone}{\mathbf{1}}
\newcommand{\bI}{\mathbf{I}}
\newcommand{\grad}{\nabla}
\newcommand{\data}{\mathcal{D}}
\newcommand\numberthis{\addtocounter{equation}{1}\tag{\theequation}}



\newcommand{\Top}[1]{\textbf{\textcolor{magenta}{Top: #1}} }
\newcommand{\yw}[1]{\textbf{\textcolor{blue}{Yingheng: #1}} }
% \DeclareMathOperator*{\argmax}{arg\,max}
% \DeclareMathOperator*{\argmin}{arg\,min}
% \newcommand{\bs}{\mathbf}

\usepackage{stackengine} 
\newcommand\oast{\stackMath\mathbin{\stackinset{c}{0ex}{c}{0ex}{\ast}{\bigcirc}}}

\let\oldtheta\theta
\renewcommand{\theta}{{\boldsymbol \oldtheta}}

\let\oldpsi\psi
\renewcommand{\psi}{{\boldsymbol \oldpsi}}

\let\oldphi\phi
\renewcommand{\phi}{{\boldsymbol \oldphi}}

\let\oldtau\tau
\renewcommand{\tau}{{\boldsymbol \oldtau}}

\let\oldomega\omega
\renewcommand{\omega}{{\boldsymbol \oldomega}}

\let\oldalpha\alpha
\renewcommand{\alpha}{{\boldsymbol \oldalpha}}

\let\oldepsilon\epsilon
\renewcommand{\epsilon}{{\boldsymbol \oldepsilon}}


\newcommand{\fix}{\marginpar{FIX}}
\newcommand{\new}{\marginpar{NEW}}

%
% These are are recommended to typeset listings but not required. See the subsubsection on listing. Remove this block if you don't have listings in your paper.
\usepackage{newfloat}
\usepackage{listings}
\DeclareCaptionStyle{ruled}{labelfont=normalfont,labelsep=colon,strut=off} % DO NOT CHANGE THIS
\lstset{%
	basicstyle={\footnotesize\ttfamily},% footnotesize acceptable for monospace
	numbers=left,numberstyle=\footnotesize,xleftmargin=2em,% show line numbers, remove this entire line if you don't want the numbers.
	aboveskip=0pt,belowskip=0pt,%
	showstringspaces=false,tabsize=2,breaklines=true}
\floatstyle{ruled}
\newfloat{listing}{tb}{lst}{}
\floatname{listing}{Listing}
%
% Keep the \pdfinfo as shown here. There's no need
% for you to add the /Title and /Author tags.
\pdfinfo{
/TemplateVersion (2025.1)
}

% DISALLOWED PACKAGES
% \usepackage{authblk} -- This package is specifically forbidden
% \usepackage{balance} -- This package is specifically forbidden
% \usepackage{color (if used in text)
% \usepackage{CJK} -- This package is specifically forbidden
% \usepackage{float} -- This package is specifically forbidden
% \usepackage{flushend} -- This package is specifically forbidden
% \usepackage{fontenc} -- This package is specifically forbidden
% \usepackage{fullpage} -- This package is specifically forbidden
% \usepackage{geometry} -- This package is specifically forbidden
% \usepackage{grffile} -- This package is specifically forbidden
% \usepackage{hyperref} -- This package is specifically forbidden
% \usepackage{navigator} -- This package is specifically forbidden
% (or any other package that embeds links such as navigator or hyperref)
% \indentfirst} -- This package is specifically forbidden
% \layout} -- This package is specifically forbidden
% \multicol} -- This package is specifically forbidden
% \nameref} -- This package is specifically forbidden
% \usepackage{savetrees} -- This package is specifically forbidden
% \usepackage{setspace} -- This package is specifically forbidden
% \usepackage{stfloats} -- This package is specifically forbidden
% \usepackage{tabu} -- This package is specifically forbidden
% \usepackage{titlesec} -- This package is specifically forbidden
% \usepackage{tocbibind} -- This package is specifically forbidden
% \usepackage{ulem} -- This package is specifically forbidden
% \usepackage{wrapfig} -- This package is specifically forbidden
% DISALLOWED COMMANDS
% \nocopyright -- Your paper will not be published if you use this command
% \addtolength -- This command may not be used
% \balance -- This command may not be used
% \baselinestretch -- Your paper will not be published if you use this command
% \clearpage -- No page breaks of any kind may be used for the final version of your paper
% \columnsep -- This command may not be used
% \newpage -- No page breaks of any kind may be used for the final version of your paper
% \pagebreak -- No page breaks of any kind may be used for the final version of your paperr
% \pagestyle -- This command may not be used
% \tiny -- This is not an acceptable font size.
% \vspace{- -- No negative value may be used in proximity of a caption, figure, table, section, subsection, subsubsection, or reference
% \vskip{- -- No negative value may be used to alter spacing above or below a caption, figure, table, section, subsection, subsubsection, or reference

\setcounter{secnumdepth}{2} %May be changed to 1 or 2 if section numbers are desired.

% The file aaai25.sty is the style file for AAAI Press
% proceedings, working notes, and technical reports.
%

% Title

% Your title must be in mixed case, not sentence case.
% That means all verbs (including short verbs like be, is, using,and go),
% nouns, adverbs, adjectives should be capitalized, including both words in hyphenated terms, while
% articles, conjunctions, and prepositions are lower case unless they
% directly follow a colon or long dash
\title{Denoising Diffusion Variational Inference: \\ Diffusion Models as Expressive Variational Posteriors}
\author{
    %Authors
    % All authors must be in the same font size and format.
    Wasu Top Piriyakulkij\equalcontrib\textsuperscript{\rm 1},
    Yingheng Wang\equalcontrib\textsuperscript{\rm 1},
    Volodymyr Kuleshov\textsuperscript{\rm 1,2}
}
\affiliations{
    %Afiliations
    \textsuperscript{\rm 1}Department of Computer Science, Cornell University\\
    \textsuperscript{\rm 2}The Jacobs Technion-Cornell Institute, Cornell Tech\\
    % If you have multiple authors and multiple affiliations
    % use superscripts in text and roman font to identify them.
    % For example,
    % Sunil Issar\textsuperscript{\rm 2},
    % J. Scott Penberthy\textsuperscript{\rm 3},
    % George Ferguson\textsuperscript{\rm 4},
    % Hans Guesgen\textsuperscript{\rm 5}
    % Note that the comma should be placed after the superscript

    % email address must be in roman text type, not monospace or sans serif
    \{wp237, yw2349, kuleshov\}@cornell.edu
%
% See more examples next
}

%Example, Single Author, ->> remove \iffalse,\fi and place them surrounding AAAI title to use it
\iffalse
\title{My Publication Title --- Single Author}
\author {
    Author Name
}
\affiliations{
    Affiliation\\
    Affiliation Line 2\\
    name@example.com
}
\fi

\iffalse
%Example, Multiple Authors, ->> remove \iffalse,\fi and place them surrounding AAAI title to use it
\title{Denoising Diffusion Variational Inference: \\ Diffusion Models as Expressive Variational Posteriors}
\author {
    % Authors
    First Author Name\textsuperscript{\rm 1},
    Second Author Name\textsuperscript{\rm 2},
    Third Author Name\textsuperscript{\rm 1}
}
\affiliations {
    % Affiliations
    \textsuperscript{\rm 1}Affiliation 1\\
    \textsuperscript{\rm 2}Affiliation 2\\
    firstAuthor@affiliation1.com, secondAuthor@affilation2.com, thirdAuthor@affiliation1.com
}
\fi


% REMOVE THIS: bibentry
% This is only needed to show inline citations in the guidelines document. You should not need it and can safely delete it.
% \usepackage{bibentry}
% END REMOVE bibentry

\begin{document}

\maketitle

\begin{abstract}
We propose denoising diffusion variational inference (DDVI), a black-box variational inference algorithm for latent variable models which relies on diffusion models
as flexible approximate posteriors. 
Specifically, our method introduces
an expressive class of diffusion-based variational posteriors that perform iterative refinement in latent space; we train these posteriors with a novel regularized evidence lower bound (ELBO) on the marginal likelihood inspired by the wake-sleep algorithm. 
Our method is easy to implement (it fits a regularized extension of the ELBO), is compatible with black-box variational inference, and outperforms alternative classes of approximate posteriors based on normalizing flows or adversarial networks.
We find that DDVI improves inference and learning in deep latent variable models across common benchmarks as well as on a motivating task in biology---inferring latent ancestry from human genomes---where it outperforms strong baselines on 1000 Genomes dataset.
\end{abstract}

\section{Introduction}

% % Latent variables...
% % Latent variable modeling is a powerful tool for representation learning, dimensionality reduction, and unsupervised learning.
% Latent variables form a powerful tool for learning representations across many application areas of machine learning [CITE], especially in scientific applications such as biological data analysis [CITE].
% % , performing dimensionality reduction, and solving problems across unsupervised learning [CITATIONS].
% Latent variable methods often rely on variational inference to fit an approximate model of the posterior distribution---the expressivity of this distribution has a significant impact on the performance of variational inference [CITE]. This fact motivates research that leverages modern generative models---including normalizing flows [CITE] and generative adversarial networks [CITE]---to represent approximate posteriors.

% Latent variables are a powerful tool for learning representations in both machine learning \citep{kingma2013auto} and its applications, including in fields of science such as biology \citep{battey2021visualizing}.

We are interested in amortized black-box variational inference 
problems of the form
\begin{align}
    \log\; p_\theta(\bx) \geq &\max_\phi \;\; \Exp_{q_\phi(\bz|\bx)} \left[ \log \; p_\theta(\bx, \bz) - \log\; q_\phi(\bz|\bx) \right] \nonumber \\ 
    &:= \max_\phi \;\; \text{ELBO}(\bx, \theta, \phi), 
\end{align}
in which we approximate the marginal likelihood $\log\; p_\theta(\bx)$ of a latent variable model $p_\theta(\bx, \bz)$ with an evidence lower bound $\text{ELBO}(\bx, \theta, \phi)$ that is a function of an approximate variational posterior $q_\phi (\bz|\bx)$.
We assume that $p_\theta$ factorizes as $p_\theta(\bx|\bz)p_\theta(\bz)$ and admits efficient sampling: examples of such $p_\theta$ include Bayesian networks, topic models \cite{blei2003latent}, variational autoencoders (VAEs), and broad classes of $p_\theta$ defined via modern probabilistic programming frameworks \cite{gordon2014probabilistic}.

Maximizing $\text{ELBO}(\bx, \theta, \phi)$ over $\phi$ yields a variational posterior $q_\phi(\bz|\bx)$ that approximates $p_\theta(\bz|\bx)$ as well as a tight bound on $\log\; p_\theta(\bx)$ that serves as a learning objective for $p_\theta$.
The approximation gap $\log\; p_\theta(\bx) - \max_\phi \text{ELBO}(\bx, \theta, \phi)$ equals precisely $\min_\phi \text{KL}(q_\phi(\bz|\bx) || p_\theta(\bz|\bx))$, which motivates the design of expressive classes of posteriors $q_\phi(\bz|\bx)$ that reduce this gap.
Recent efforts %towards this goal 
leverage modern generative models---including normalizing flows \citep{rezende2015variational,kingma2016improved} and generative adversarial networks \citep{goodfellow2014generative,makhzani2015adversarial}---as expressive model families for $q_\phi$ that tighten the ELBO.

This work seeks to further improve variational inference via expressive posteriors based on diffusion models
\citep{ho2020denoising, song2020score}.
Diffusion methods have become 
the de-facto standard for high-quality image synthesis \citep{rombach2022high,gokaslan2024commoncanvas}.
Here, we use diffusion in latent space to parameterize $q_\phi(\bz|\bx)$.
We train this distribution with a denoising diffusion-like objective that does not involve adversarial training \citep{makhzani2015adversarial} or constrained invertible normalizing flow architectures \citep{kingma2016improved}. Samples from $q_\phi(\bz|\bx)$ are obtained via iterative refinement of $\bz$, starting from a Gaussian distribution, and gradually forming one that is multi-modal and complex.


Our work expands upon
existing diffusion-based approximate inference methods \cite{berner2022optimal, zhang2021path, vargas2023denoising, zhang2023diffusion, richter2023improved, sendera2024diffusion, akhound2024iterated} that focus on the task of drawing samples from unnormalized distributions $\tilde p(\bz)$ and estimating the partition function $Z = \int_\bz \tilde p(\bz) d\bz$. 
While these methods are applicable in our setting---we set the unnormalized $\tilde p(\bz)$ to $p_\theta(\bx,\bz)$ such that $Z = p_\theta(\bx)$---they do not make use of characteristics of $p_\theta$ that are common in many types of models (VAEs, Bayes networks, etc.), namely the factorization $p_\theta(\bx|\bz)p_\theta(\bz)$ and efficient sampling.
We find that leveraging these properties yields simpler algorithms that avoid backpropagating through a sampling process, and that are fast enough to perform learning in addition to inference.

Specifically, we propose denoising diffusion variational inference (DDVI), an approximate inference algorithm defined by a class of approximate posterior distribution based on diffusion and a learning objective inspired by the wake-sleep algorithm \citep{hinton1995wake} that implements regularized variational inference.
We also derive extensions of our method to semi-supervised learning and clustering.

Our method is easy to implement (it fits a regularized extension of the ELBO), is compatible with black-box variational inference, and outperforms alternative classes of approximate posteriors based on normalizing flows or adversarial networks.

We evaluate DDVI on synthetic benchmarks and on a real problem in biological data analysis---inferring human ancestry from genetic data. Our method outperforms strong baselines on 1000 Genomes dataset \citep{siva20081000} and learns a low-dimensional latent space that preserves biologically meaningful structure \citep{haghverdi2015diffusion}.

\paragraph{Contributions.} In summary, this work introduces denoising diffusion variational inference, an approximate inference algorithm defined by two components: 
a class of approximate posteriors $q(\bz|\bx)$ parameterized by diffusion,
and a lower bound on the marginal likelihood inspired by wake-sleep.
We complement DDVI with extensions to semi-supervised learning and clustering.
Our method is especially suited for probabilistic programming, representation learning, and dimensionality reduction, where it outperforms alternative methods based on normalizing flows and adversarial training.

\begin{figure*}[t]
    \centering
    \includegraphics[width=0.9\textwidth]{imgs/diffvae_figure_fixed.pdf}
    %\includegraphics[scale=1.5]{spiral}

    \caption{Denoising diffusion variational inference in a VAE. Between the encoder and decoder, we have a diffusion model to map a simple distribution into a complex distribution over latents.}
    \label{fig:diffvae_main}
\end{figure*}

% \begin{itemize}
% \item Visualization is an important task. Many scientific tasks require visualizing data in a low-dimensional space. One approach is via latent variable models.
% \item In order to be effective latent variable models need to be sufficiently expressive to represent the posterior distribution $p(z|x)$ of latent variables. This is usually done via variational infernece.
% \item People have tried to improve variational inference using more expressive models (such as normalizing flows), or using more suitable objective (adversarial autoencoders, MMD-VAE). However, these approaches can sometimes be unstable or computationally intensive. \Top{And for normalizing flow, approximating the KL term for non-gaussian prior requires careful weighting between $E_{q(z|x)} [\log q(z|x)]$ and $E_{q(z|x)} [\log p(z)]$ for each different prior} \yw{How could we justify "computationally intensive"? Diffusion models are also computationally intensive in some sense, especially when we incorporate sampling from reverse diffusion distribution for training.}
% \item Recently, diffusion models have emerged as a method that combines high sample quality with a stable training objective.
% \item This paper explores the use of diffusion models in latent space as variational posteriors. We expect diffusion models to render complex approximate latent variable posteriors similarly to how they can render complex images.
% \item Applying DMs naively within the ELBO is challenging: it requires backpropagating through potentially thousands of steps of Langevin MCMC sampling. 
% \item Instead, our approach develops an alternative objective derived from applying the ELBO followed by an alternative lower bound derived from the wake-sleep algorithm. This yields an end-to-end differentiable objective combining the benefits of variational inference and wake sleep. We derive additional extensions of this objective that are compatible with semi-supervised learning, learnable priors, and other attractive features. \Top{It should be noted that the objective of wake-sleep algorithm is not actually a lower bound. It optimizes an objective that is almost similar to a lower bound but using the wrong, reverse KL divergence.}
% \item While our method does not yield density estimates that are similar to ones derived from methods optimizing the ELBO, we are able to fit highly complex posteriors better than existing alternatives based on adversarial training or normalizing flows.
% \item We show that this method is useful for a range of visualization tasks, on both synthetic dataset, and a real dataset in genomics.
% \end{itemize}

\section{Background}

% \subsection{Latent Variable Models and Variational Inference}

\paragraph{Deep Latent Variable Models}

Latent variable models (LVMs)  $p_\theta(\bx, \bz)$ are usually fit by optimizing the evidence lower bound (ELBO) 
\begin{equation}
\log p_\theta(\bx) \geq \mathbb{E}_{q_\phi(\bz|\bx)} [\log p_\theta(\bx|\bz)] - \KL(q_\phi(\bz|\bx)||p_\theta(\bz)), \nonumber
\end{equation}
which serves as a tractable surrogate for the marginal log-likelihood (MLL). The gap between the MLL and the ELBO equals precisely $\KL(q_\phi(\bz|\bx)||p_\theta(\bz|\bx))$---thus, a more expressive $q_\phi(\bz|\bx)$ may better fit the true posterior and induce a tighter ELBO \citep{kingma2013auto}.

Expressive variational posteriors can be formed by choosing more expressive model families---including auxiliary variable methods \citep{maaloe2016auxiliary}, MCMC-based methods \citep{salimans2015mcmc}, normalizing flows \citep{rezende2015variational}---or improved learning objectives---e.g., adversarial or sample-based losses \citep{makhzani2015adversarial,zhao2017infovae,si2022autoregressive,pmlr-v202-si23a}.
% Multiple strategies proposed for approximating the posteriors have gained prominence, including normalizing flows which transform simple distributions into complex ones, Markov Chain Monte Carlo (MCMC) methods for sampling-based approximations, and the use of auxiliary variables to aid in model expressivity. Among these, AAE \citep{makhzani2015adversarial} and IAF-VAE \citep{kingma2016improved} stand out. An alternative line of work focuses on modifying the VAE objective itself \citep{zhao2017infovae, makhzani2015adversarial}. 

The wake-sleep algorithm \citep{hinton1995wake} optimizes an alternative objective
\begin{equation}
\mathbb{E}_{q_\phi(\bz|\bx)} [\log p_\theta(\bx|\bz)] - \KL(p_\theta(\bz|\bx)||q_\phi(\bz|\bx)), \nonumber
\end{equation}
in which the KL divergence term is reversed. The learning procedure for wake-sleep involves alternating between "wake" phases where the recognition model is updated and "sleep" phases where the generative model is refined.

% \begin{itemize}
%     \item First mention ELBO for an LVM $p(\bx,\bz)$.
%     \item Then explain why better posterior yields a tighter bound
%     \item Then explain existing methods for fitting this posterior: especially make sure to call out AAE and VAE-IAF. Talk about various model classes for approximate posteriors: normalizing flows, MCMC, auxilairy variables. Talk about another set of approaches that modify the objective: MMD-VAE, AAE, InfoVAE
%     \item then mention wake sleep as another algorithm that optimizes a bound similar to the ELBO, except that the regularization term changes and the KL gets flipped (give the bound explicitly). define the objective and describe in one sentence the learning procedure
% \end{itemize}

\paragraph{Denoising Diffusion Models}

A diffusion model is defined via a user-specified noising process $q$ that maps data $\bx_0$ into a sequence of $T$ variables $\by_{1:T} = \by_1, ..., \by_T$ that represent increasing levels of corruption to $\bx_0$. We obtain $\by_{1:T}$ by applying a Markov chain
$q(\by_{1:T} | \bx_0) = \prod_{t=1}^T q(\by_t | \by_{t-1})$, where we define $\by_0 = \bx_0$ for convenience.
% {\footnotesize
% \begin{equation}
% q(\bx_{1:T} \mid \bx_0) = \prod_{t=1}^T q(\bx_t \mid \bx_{t-1}).
% \end{equation}
% }
When $\bx_0$ is a continuous vector, a standard choice of transition kernel is $q(\bx_t \mid \bx_{t-1}) = \mathcal{N}(\by_t; \sqrt{\alpha_t} \by_{t-1}, \sqrt{1-\alpha_t} \mathbf{I})$, which is a Gaussian centered around a copy of $\by_{t-1}$ to which we added noise following a schedule $0<\alpha_1 < \alpha_2 < ... < \alpha_T=1$.
% As shown in \citet{ho2020denoising}, the marginal distribution of $q$ can be expressed as 
% \vspace{0.1cm}
% $$
% q(\bx_t \mid \bx_0) = \mathcal{N}(\bx_t; \sqrt{\bar \alpha_t} \bx_0, \sqrt{1-\bar \alpha_t} \mathbf{I}),
% $$
% \vspace{0.1cm}
% where $\bar\alpha_t = \prod_{s=1}^t \alpha_t$ is the cumulative product of the schedule parameters $\alpha_t$.

A diffusion model can then be represented as a latent variable distribution $p(\bx_{0}, \by_{1:T})$ % over data $\bx_0$ as well as latents $\by_{1:T} := \by_1, \by_2, \ldots, \by_T$ that represent a gradual transformation of $\bz_0$ into  random Gaussian noise $\by_T$.
% The distribution $p$ 
that factorizes as
$p(\bx_0, \by_{1:T}) = p(\by_{T}) \prod_{t=0}^{T-1} p_\theta(\by_t \mid \by_{t+1})$ (again using $\by_0$ as shorthand for $\bx_0$).
% {\footnotesize
% \begin{equation}
% p(\bx, \by_{1:T}) = p(\bx_{T}) \prod_{t=0}^{T-1} p_\theta(\bx_t \mid \bx_{t+1}).
% \end{equation}
% }
This model seeks to approximate the reverse of the forward diffusion $q$ and map noise $\by_T$ into data $\bx_0$.
% by ``undoing" a noising (or diffusion) process denoted by $q.$
% Here we use a learned denoising distribution $p_\theta,$ which we parameterize by a neural network with parameters $\theta$. 

The true reverse of the process $q$ cannot be expressed in closed form; as such, we parameterize $p_\theta$ with $\theta$
trained by maximizing the ELBO:
\vspace{-5pt}
\begin{align}
    \log p_\theta(\bx_{0}) \geq &
    \Exp_{q} \Bigr[\log p_\theta(\bx_{0} | \bx_{1})  -  \sum_{t=2}^T  \KL(q_t || p_t) \Bigr] \\ \nonumber
    &- \KL(q(\bx_T | \bx_{0}) || p(\bx_T))
\end{align}
where $q_t, p_t$ denote the distributions $q(\bx_{t-1} | \bx_{t}, \bx_0)$ and $p_\theta(\bx_{t-1} | \bx_{t})$, respectively.
% Notice that the ELBO objective fits each $p_\theta(\bx_{t-1} | \bx_{t})$ to the $t$-th step of the true reverse diffusion process $q(\bx_{t-1} | \bx_{t}, \bx_0)$.

% \Top{Should we also mention the noise prediction formula here?}

% \subsection{Visualization and Dimensionality Reduction}

% \paragraph{Visualization and Dimensionality Reduction} There exist two important types of dimensionality reduction methods, which (1) emphasize the preservation of pairwise distance structures among all data samples, e.g., PCA \citep{wold1987principal} and LDA \citep{balakrishnama1998linear}, and (2) local distances over global ones, e.g., t-SNE \citep{van2008visualizing} and UMAP \citep{mcinnes2018umap}. Latent variable models (LVMs) \citep{kingma2016improved, makhzani2015adversarial} are another class of techniques for dimensionality reduction. They represent high-dimensional data in terms of latent, or hidden, variables in a lower-dimensional space, effectively providing a compact representation of the data. 
% However, the standard LVMs often assume simple Gaussian priors and posteriors. As the complexity of real-world data increases, there is a need for a more expressive variational posterior to capture intricate data distributions and relationships, ensuring the latent space representation remains meaningful and informative.

% \begin{itemize}
%     \item Briefly introduce key types of dimensionality reduction
%     \item Explain how LVMs can be used for DR. Explain the need for a more expressive variational posterior.
% \end{itemize}

\section{Variational Inference With\\Denoising Diffusion Models}
\label{sec:diffvae}

We introduce {\em denoising diffusion variational inference (DDVI)}, which improves variational inference with diffusion-based techniques.

The goal of DDVI is to fit a latent variable model $p_\theta(\bx,\bz)$.
We assume that $p_\theta$ factorizes as $p_\theta(\bx|\bz)p_\theta(\bz)$ and admits efficient sampling: examples of such $p_\theta$ include Bayesian networks and variational autoencoders (VAEs) \cite{kingma2013auto}.

Our approach is comprised of three components: 
\begin{enumerate}
    \item A \textbf{modeling family} of approximate posteriors $q_\phi(\bz | \bx)$ based on diffusion;
    \item A \textbf{learning objective} formed by a regularized ELBO;
    \item An \textbf{optimization algorithm} inspired by wake-sleep.
\end{enumerate}
The $q_\phi(\bz | \bx)$ iteratively refines latents $\bz$, starting from a Gaussian distribution. 
The learning objective trains $q_\phi(\bz | \bx)$ to reverse a used-specified forward diffusion process.

\subsection{Modeling Family: Diffusion-Based Posteriors}

DDVI performs variational inference using a family of approximate posteriors $q_\phi(\bz|\bx) = \int_\by q_\phi(\bz| \by,\bx) q_\phi(\by|\bx) d\by$, which themselves contain latent variables $\by \in \mathcal{Y}$.
The models $q_\phi(\bz| \by,\bx), q_\phi(\by|\bx)$ must have tractable densities and support gradient-based optimization over $\phi$.

We choose the latent $\by = (\by_1, \by_2, ..., \by_T)$ to be  a vector of $T$ variables that represent progressively simplified versions of $\bz$, with $\by_T$ corresponding to a simple distribution (e.g., a Gaussian).
The model $q_\phi(\by, \bz|\bx) = q_\phi(\bz| \by_1,\bx) \prod_{t=1}^{T-1} q_\phi(\by_t|\by_{t+1},\bx)$ transforms $\by_T$ into $\bz$ via iterative refinement.
To sample from $q_\phi$, we first sample $\by_T$---this is an easier task since we can define $\by_T$ to have a simple (e.g., Gaussian) distribution---and then by sampling from the denoising model $q_\phi(\bz| \by_1,\bx) \prod_{t=1}^{T-1} q_\phi(\by_t|\by_{t+1},\bx)$. 

We define the relationship between $\by$ and $\bz$ via a forward diffusion process $r(\by|\bz,\bx) = r(\by_{1:T}|\bz,\bx) =  r(\by_1|\bz,\bx)\prod_{t=1}^{T-1} r(\by_{t+1}|\by_{t},\bx) $,
which transforms $\bz$---the latent whose intractable posterior we seek to approximate---into $\by_T$, whose posterior is easier to model (possibly conditioned on $\bx$).
Examples of $r$ include Gaussian forward diffusion processes and discrete noising processes \citep{austin2021structured}.
The model $q_\phi$ is trained to approximately reverse this forward diffusion process.

\subsection{Learning Objective: A Markovian ELBO}

The standard approach to fit auxiliary-variable generative models \citep{maaloe2016auxiliary} is to apply the ELBO twice:
\begin{align}
    \log p_\theta(\bx) 
    & \geq \log p_\theta(\bx) - \KL(q_\phi(\bz|\bx)||p_\theta(\bz|\bx)) \label{eqn:elbo2_z1} \\
    & \geq \log p_\theta(\bx) - \KL(q_\phi(\bz|\bx)||p_\theta(\bz|\bx)) \label{eqn:elbo2_y1}  \\ \nonumber
    &~~~~ - \Exp_{q_\phi(\bz|\bx)} [ \KL(q_\phi(\by|\bx,\bz)||r(\by|\bx,\bz)) ] \\
    & = ~\Exp_{q_\phi(\by,\bz|\bx)} [\log p_\theta(\bx|\bz)] \label{eqn:elbo2_y2} \\ \nonumber
    &~~~~ -\KL(q_\phi(\by,\bz|\bx)||r(\by|\bx,\bz)p(\bz))
\end{align}
In Equation (\ref{eqn:elbo2_z1}), we applied the ELBO over $\bz$, and in Equation (\ref{eqn:elbo2_y1}) we applied the ELBO over the latent $\by$ of $q$ (see Appendix J for the derivation). 
Notice that the gap between the ELBO and $\log p_\theta(\bx)$ is $\KL(q_\phi(\bz|\bx)||p_\theta(\bz|\bx)) + \Exp_{q_\phi(\bz|\bx)} [ \KL(q_\phi(\by|\bx,\bz)||r(\by|\bx,\bz)) ]$. Thus, if we correctly match $q$ and $r$, we will achieve a tight bound.

\subsubsection{Analyzing the ELBO}
% Lastly, in order to use diffusion-based encoders with our ELBO-based objective, 
Optimizing Equation (\ref{eqn:elbo2_y2}) requires tractably dealing with the prior regularization term $\mathcal{L}_\text{reg}(\bx,\theta,\phi) := -\KL(q_\phi(\by,\bz|\bx)||r(\by|\bx,\bz)p(\bz))$, which we equivalently rewrite as:
\begin{align}
\mathcal{L}_\text{reg} = \Exp_{q_\phi(\by,\bz|\bx)}[\log(r(\by|\bx,\bz)p(\bz))] + H(q). \label{eqn:reg_loss}
\end{align}
We can expand the first term by leveraging the Markov structure of  $r, q$ to rewrite $\mathcal{L}_\text{reg}$ as the likelihood of samples from the reverse diffusion process $q$ under the forward process $r$.
\begin{align}
\mathcal{L}_\text{reg} = \sum_{t=1}^T \Exp_{q}[\log(r(\by_t|\by_{t-1},\bx)] + \Exp_{q} [\log p(\bz)] + H(q), \nonumber
\end{align}
where $\by_0 := \bz$. We refer to optimizing the Markovian ELBO as \textbf{unregularized DDVI}, the first instance of our method. 
% This approach tightens the ELBO via expressive $q$'s with latent variables; 
The noise process $r$ defines prior regularization terms for each $\by_t$. This provides extra supervision for learning $q$ in the form of trajectories from latents $\by_T$ to $\by_1$; this extra supervision helps $q$ fit complex non-Gaussian posteriors.
% The noise process $r$ helps fit expressive $q$'s with latent variables by providing supervision via trajectories from latents $\by_T$ that are easy to model to latents $\by_1$ that are complex.
% The $r$ provides extra supervision for learning $q$ and helps it fit complex priors.
% Recall that $
% r(\by_{1:T}|\bz,\bx) =  \prod_{t=1}^T r(\by_t|\by_{t-1},\bx) 
% $,
% where $\by_0 = \bz$ and
% $
% q_\phi(\by_{0:T} | \bx) = q_\phi(\by_T | \bx) \prod_{t=1}^T q_\phi(\by_{t-1}|\by_{t},\bx)
% $.
% We may use the Markov structure in $q, r$ to rewrite the first term in (\ref{eqn:reg_loss}) as a sum of $T$ terms, one per Markov step. The derivation is identical to that used to obtain the ELBO of a diffusion model \citep{sohl2015deep}, and yields an expression of the same form:
% % In order to fit $q_\phi$,
% % we form a lower bound $\Lc_\text{diff}(\bx,\phi)$ on the sleep term $\Exp_{p(\bx,\bz)} \log {q_\phi(\bz|\bx)}$ in Equation \ref{eqn:sleep1b}.
% % This bound is identical to the ELBO of a diffusion model, and has the same derivation \citep{sohl2015deep}.
% \begin{align}
% \Lc_\text{diff} = &~\Exp_{r} \left[\log q_\phi(\bz | \by_{1}, \bx)  -  \sum_{t=2}^T  \KL(r_t || q_t) \right] \\ \nonumber
% &- \KL(r(\by_T | \bz) || q_\phi(\by|\bx)).
% \end{align}
% where $r_t,q_t$ denote the distributions $r(\by_{t-1} | \by_{t}, \by_0, \bx)$ and $q_\phi(\by_{t-1} | \by_{t}, \bx)$.

The term $H(q) = - \sum_{t=1}^{T+1} \mathbb{E}_q [\log q_\phi(\by_{t-1}|\by_t,\bx)]$ denotes the entropy. For example, when each term $q_\phi(\by_{t-1}|\by_t,\bx)$ is Gaussian, it is computed as
\begin{align}
    H(q) = 
    % &- \sum_{t=1}^{T+1} \mathbb{E}_q [\log q_\phi(\by_{t-1}|\by_t,\bx)] \\ \nonumber =
    & \sum_{t=1}^{T+1} \mathbb{E}_{q} \left [ \frac{d}{2} \left( 1 + \log(2\pi) \right) + \frac{1}{2} \log \left| \Sigma_\phi(\by_t, \bx) \right| \right ] \nonumber
\end{align}
where $d$ is the dimension of $\by$ and we use the notation $\by_{T+1} = \bx$. 
% The right-hand term can be approximated using Monte Carlo; 
It is also common to leave the variance $\Sigma_\phi$ fixed, in which case $H(q)$ is a constant.

\subsection{Refining the Objective: A Regularized ELBO}

Notice that optimizing $\mathcal{L}_\text{reg}$
% $\KL(q_\phi(\by,\bz|\bx)||r(\by|\bz,\bx)p(\bz))$ in Equation \ref{eqn:elbo2_y2} 
involves sampling from the approximate reverse process $q_\phi(\by,\bz|\bx)$ to match the true reverse process $r(\by | \bz, \bx)$: this is the opposite of diffusion training, where we would sample from $r$ to fit $q$. This type of on-policy learning of $q$ has been studied in the context of approximate inference \citep{zhang2021path}; however, it requires backpropagating through $T$ samples, which may hamper training, and it optimizes a mode-covering divergence that may struggle to fit complex $p(\bz)$. 
% We highlight challenges with on-policy optimization  in Table \ref{tab:unsupervised_w_failed} in Appendix.
% We have found the standard ELBO to insufficient to learn a good $q_\phi$ that reverses the noising process $r$, 
% We found that on-policy learning of $q$ often struggles to learn a good $q_\phi$, as illustrated in Table \ref{tab:unsupervised_w_failed}. % in Appendix.

\subsubsection{Adding Wake-Sleep Regularization to the ELBO}

We propose to further improve the ELBO via off-policy diffusion-like training. Our new objective is the ELBO in Equation (\ref{eqn:elbo2_y2}) augmented with a regularizer $\Lc_\text{sleep}(\phi)$.
\begin{align} \label{eqn:objective}
\log p_\theta(\bx) \geq &
 \underbrace{\Exp_{q_\phi(\by,\bz|\bx)} [\log p_\theta(\bx|\bz)]}_\text{wake / recons. term $\Lc_\text{rec}(\bx,\theta,\phi)$} \\ \nonumber & \underbrace{-\KL(q_\phi(\by,\bz|\bx)||r(\by|\bx,\bz)p(\bz))}_\text{prior regularization term $\mathcal{L}_\text{reg}(\bx,\theta,\phi)$} \\ \nonumber & \underbrace{-\Exp_{p_\theta(\bx)} [\KL(p_\theta(\bz|\bx)||q_\phi(\bz|\bx))]}_\text{sleep term $\Lc_\text{sleep}(\phi)$}
\end{align}
The optimization of the regularizer $\Lc_\text{sleep}(\phi)$ is similar to the sleep phase of wake-sleep, and closely resembles diffusion model training (see below). 
As in wake-sleep, $\Lc_\text{sleep}(\phi)$ is optimized over $\phi$ only, the $\bx$ are sampled from $p$.

\subsubsection{From Wake-Sleep to Diffusion Regularization}

Computing $\Lc_\text{sleep}(\phi)$ still involves intractable distributions $p_\theta(\bz|\bx), q_\phi(\bz|\bx)$.
To optimize $\Lc_\text{sleep}(\phi)$, we introduce another lower bound $\Lc_\text{diff}(\phi)$, which we call the denoising diffusion loss (for reasons that will become apparent shortly):
\begin{align}
\label{eqn:sleep1b}
\Lc_\text{sleep}(\phi) 
& = - \Exp_{p_\theta(\bx)} [\KL(p_\theta(\bz|\bx)||q_\phi(\bz|\bx))] \\ \nonumber
& = \Exp_{p_\theta(\bx,\bz)} [\log {q_\phi(\bz|\bx)}] + \bar H(p_\theta) \\ \nonumber
& \geq \Exp_{p_\theta(\bx,\bz)} \left[ \Exp_{r} [\log \frac{ q_\phi(\by,\bz|\bx)} {r(\by|\bz,\bx) }] \right] + \bar H(p_\theta) \\ \nonumber
& = \Lc_\text{diff}(\phi)
\end{align}
In Equation (\ref{eqn:sleep1b}), we applied the ELBO with $r(\by|\bz,\bx)$ playing the role of the variational posterior over the latent $\by$ in $q_\phi$; $\bar H(p_\theta)$ is the expected conditional entropy of $p_\theta(\bz|\bx)$, a constant that does not depend on $\phi$.

% \subsubsection{Simplifying the Diffusion Term}

% Lastly, w
We can further simplify $\Lc_\text{diff}(\phi)$ by leveraging the Markov structure of the forward and reverse processes $r, q$.
Recall that each $\by = (\by_1, \by_2, ..., \by_T)$ can be a vector of $T$ latents, which we also denote as $\by_{1:T}$,
and that
$
r(\by_{1:T}|\bz,\bx) =  \prod_{t=1}^T r(\by_t|\by_{t-1},\bx) 
$,
where $\by_0 = \bz$ % and the $\by_{1:T}$ are increasingly noised versions of $\by_0$.
and also
% $q_\phi(\by, \bz | \bx)$ is an approximate reverse diffusion process
$
q_\phi(\by, \bz | \bx) = q_\phi(\by_{0:T} | \bx) = q_\phi(\by_T | \bx) \prod_{t=1}^T q_\phi(\by_{t-1}|\by_{t},\bx).
$
% that is trained to reverse $r$.

We may use the Markov structure in $q, r$ to rewrite $\Lc_\text{diff}(\phi)$ as a sum of $T$ terms, one per Markov step. The derivation is identical to that used to obtain the ELBO of a diffusion model \citep{sohl2015deep}, and yields an expression of the same form:
% In order to fit $q_\phi$,
% we form a lower bound $\Lc_\text{diff}(\bx,\phi)$ on the sleep term $\Exp_{p(\bx,\bz)} \log {q_\phi(\bz|\bx)}$ in Equation \ref{eqn:sleep1b}.
% This bound is identical to the ELBO of a diffusion model, and has the same derivation \citep{sohl2015deep}.
\begin{align}
\Lc_\text{diff}(\phi) = &~\Exp_{r} \left[\log q_\phi(\bz | \by_{1}, \bx)  -  \sum_{t=2}^T  \KL(r_t || q_t) \right] \\ \nonumber
&- \KL(r(\by_T | \bz, \bx) || q_\phi(\by_T | \bx)) + \bar{H}(p_\theta).
\end{align}
where $r_t,q_t$ denote the distributions $r(\by_{t-1} | \by_{t}, \by_0, \bx)$ and $q_\phi(\by_{t-1} | \by_{t}, \bx)$ (see Appendix K for the derivation).

\subsubsection{Regularized DDVI Objective}

We define the full DDVI objective $\Lc_\text{ddvi}$ to be the sum of the aforementioned terms:
\begin{align}
    \Lc_\text{ddvi}(\bx, \theta, \phi) = \Lc_\text{rec}(\bx, \theta, \phi) + \Lc_\text{reg}(\bx, \theta, \phi) + \Lc_\text{diff}(\phi) \nonumber
\end{align}
Terms $\Lc_\text{reg}$ and $\Lc_\text{diff}$ may be weighted by hyper-parameters $\beta_\text{reg}, \beta_\text{diff} > 0$, as in the $\beta$-VAE framework. In our experiments, $\beta_\text{reg} = \beta_\text{diff} = 1$ unless otherwise specified.
Note that since $\Lc_\text{diff} \leq \Lc_\text{sleep} \leq 0$, $\Lc(\bx, \theta, \phi)$ is a valid lower bound on $\log p_\theta(\bx)$ that is tight when $q_\phi(\bz|\bx) = p_\theta(\bz|\bx)$.



\subsection{Optimization: Extending Wake-Sleep}\label{sec:latent-wake-sleep}

% \paragraph{Wake-Sleep Optimization}
We may optimize $\Lc_\text{ddvi}(\bx, \theta, \phi)$ using gradient descent by alternating between ELBO optimization and taking sleep steps (see Appendix A for full details):
\begin{enumerate}
    \item Sample $\bx$ from data, take gradient step on $\theta, \phi$ optimizing $\Lc_\text{rec}(\bx, \theta, \phi) + \Lc_\text{reg}(\bx, \theta, \phi)$ (the ``wake" step);
    \item Sample $\bz, \bx$ from $p_\theta$ and take a gradient step on $\phi$ optimizing $\Lc_\text{diff}(\phi)$ (the ``sleep" step).
\end{enumerate}
Again, terms may be weighted by $\beta_\text{reg}, \beta_\text{diff} > 0$.
Note that by maximizing $\Lc_\text{diff}(\phi)$, we fit $q_\phi(\bz|\bx)$ to $p_\theta(\bz|\bx)$ via the forward KL divergence; similarly, by optimizing $\Lc_\text{rec} + \Lc_\text{reg}$ (the ELBO), we fit $q_\phi(\bz|\bx)$ to $p_\theta(\bz|\bx)$ via the reverse KL divergence. Thus, optimizing $\Lc(\bx, \theta, \phi)$ encourages $q_\phi(\bz|\bx)$ to approximate $p_\theta(\bz|\bx)$, and when the two are equal, the bound $\Lc_\text{ddvi}$ on $\log p_\theta(\bx)$ is tight.

\paragraph{Simplified Wake-Sleep}
We also consider a light-weight algorithm, in which $r(\by|\bz)$ and $q_\phi(\by, \bz)$ do not depend on $\bx$. This scenario admits the following optimization procedure:
\begin{enumerate}
    \item Sample $\bx$ from data and compute gradient on $\theta, \phi$ optimizing $\Lc_\text{rec}(\bx, \theta, \phi) + \Lc_\text{reg}(\bx, \theta, \phi)$.
    \item Sample $\bz$ from $p(\bz)$ and compute gradient on $\phi$ optimizing $ \Lc_\text{diff}(\phi)$; take step on weighted sum of both gradients.
\end{enumerate}
In this case, $\Lc_\text{diff}$ requires only sampling from $p(\bz)$, and the entire loss $\Lc_\text{ddvi}$ can be optimized end-to-end using gradient descent. This algorithm is simpler (there is no separate sleep phase); however, $q_\phi(\bz|\bx)$ may not perfectly approximate $p_\theta(\bz|\bx)$ when $r(\by|\bz)$ and $q_\phi(\bz|\by)$ do not depend on $\bx$, hence $\Lc$ may no longer be a tight bound.





\subsubsection{Practical Considerations}

A common type of noising process compatible with this bound when $\bz$ is continuous is Gaussian diffusion, where we define $r(\by_t | \by_{t-1}) = \mathcal{N}(\by_t; \sqrt{1 - \alpha_t} \by_{t-1}, \alpha_t \bI)$ for a suitable schedule $(\alpha_t)_{t=1}^T$. We then adopt the parameterization $q_\phi(\by_{t-1} | \by_t, \bx) = \mathcal{N}(\by_{t-1}; \mu_\phi (\by_t, \bx, t), \Sigma_\phi (\by_t, \bx, t))$. 
It is then common to parameterize $q_\phi$ with a noise prediction network $\epsilon_\phi$ \citep{ho2020denoising}; the sum of KL divergences can be approximated by $\mathbb{E}_{t, \epsilon_t \sim r(\by_0, t)}||\epsilon_t - \epsilon_\phi(\sqrt{\bar{\alpha_t}} \by_0 + \sqrt{1-\bar{\alpha_t}} \epsilon_t, \bx, t)||^2$. Other extensions include discrete denoising processes \citep{austin2021structured,sahoo2024simple,schiff2024simple}. In the wake-sleep setting, we know both endpoints $\by_T \sim q(\cdot|\bx)$ and $\by_0 = \bz$ of the diffusion process, opening the possibility for applying optimal transport techniques \citep{cuturi2013sinkhorn,de2021diffusion}.



\section{Extensions}

\subsection{Semi-Supervised Learning}

Following \citet{makhzani2015adversarial}, we extend our algorithm to the semi-supervised learning setting where some data points have labels denoted by $l$. We assume the user provides a model of the form $p_\theta(\bx, \by, \bz, l) = p_\theta(\bx|\bz, l)r(\by|\bz, l)p_\theta(\bz|l)p(l)$; we set the variational distributions to $q_\phi(\bz|\bx, \by, l), q_\phi(\by|\bx), q_\phi(l|\bx)$. In this setting, we consider two cases, depending on whether the label is observed \citep{kingma2014semi}. We extend Equation (\ref{eqn:objective}) to incorporate the label $l$ corresponding to a data point as follows:
\begin{align}
    \Lc_\text{semi} =
    &~\Exp_{q_\phi(\by,\bz|\bx, l)} [\log p_\theta(\bx|\bz, l)] \\ \nonumber 
    &-\KL(q_\phi(\by,\bz|\bx, l)||p_\theta(\by,\bz|l)) \\ \nonumber
    &- \Exp_{p_\theta(\bx)} \left [ \KL(p_\theta(\bz|\bx, l)||q_\phi(\bz|\bx, l)) \right ]
\end{align}
When the label $c$ cannot be observed, we treat it as a latent variable and modify the learning objective $\Uc_\text{semi} 
    = \sum_c q_\phi(l|\bx)\Lc_\text{semi}(\bx, l, \btheta, \phi) + \KL(q_\phi(l|\bx)||p(l))$. Therefore, we can conclude a marginal likelihood on our dataset as follows:
$
    \tilde{\Lc}_\text{semi} = \sum_{(\bx, l) \in L}\Lc_\text{semi}(\bx, l, \btheta, \bphi) + \sum_{\bx \in U}\Uc_\text{semi}(\bx, \btheta, \bphi).
$
where $L$ and $U$ are the sets of data with and without labels, respectively.

We also want to guarantee that all model parameters can be learned in all cases, including $q_\phi(l|\bx)$, such that this posterior can be applied as a classifier during inference. Thus, we combine the marginal likelihood with a classification loss to form an extended learning objective:
$
    \tilde{\Lc}_{\text{semi}_\alpha} = \tilde{\Lc}_\text{semi} + \alpha \cdot \Exp_{\tilde{p}(\bx, l)} \left[ - \log q_\phi (l|\bx) \right]
$

\subsection{Clustering}
We have further extended our algorithm to encompass the clustering paradigm. We propose two distinct strategies. In the first approach, we simply formulate a model in which $p_\btheta(\bz)$ is a mixture of desired priors. The means of these priors are characterized by $\btheta$. From these means, cluster membership, denoted as $\bc$ can be deduced. This approach requires no alteration to the existing learning objective.

Alternatively, the second method retains the original prior, but introduces an additional latent cluster variable $\bc$ where $\sum_i c_i = 1$. Thus, the model can be specified as $p_\theta(\bx, \by, \bz, \bc) = p_\theta(\bx|\bz, \bc)r(\by|\bz)p_\theta(\bz)p(\bc)$ with $p(\bc) = Dir(\bepsilon)$. Consequently, the variational distributions become $q_\phi(\bz|\by, \bc, \bx), q_\phi(\by, \bc|\bx)$. This yields the objective:
\begin{align}
    \Lc_\text{clus} (\bx) =
    &~\Exp_{q_\phi(\by,\bz, \bc|\bx)} [\log p_\theta(\bx|\bz, \bc)] \\ \nonumber
     &- \KL(q_\phi(\by,\bz,\bc|\bx)||p_\theta(\by,\bz,\bc)) \\ \nonumber &- \Exp_{p_\theta(\bx)} \left [ \KL(p_\theta(\bz | \bx)||q_\phi(\bz|\bx)) \right ]
\end{align}
Expectations over small numbers of classes $\bc$ are done analytically; larger $\bc$ require backpropagating through discrete sampling \citep{jang2016categorical,sahoo2023backpropagation}.
\begin{figure*}[htbp!]
    \centering
    % \includegraphics[width=0.95\textwidth]{imgs/unsupervised_fig_cropped_add_pis.pdf}
    \includegraphics[width=0.95\textwidth]{imgs/unsupervised_no_font_cropped.pdf}
    % \vspace{-5pt}
    
    %\includegraphics[scale=1.5]{spiral}
    \caption{Unsupervised visualization on MNIST using three priors (pinwheel, swiss roll, and square). Each color indicates a class.}
    % \vspace{-10pt}
    \label{fig:unsupervised}
\end{figure*}

\begin{table*}[ht]
\centering
\adjustbox{max width=0.95\textwidth}{
\centering
\setlength\tabcolsep{4pt} % control the horizontal spacing between columns
\begin{tabular}{c@{\hspace{5pt}}c@{\hspace{5pt}}c@{\hspace{5pt}}c@{\hspace{5pt}}c@{\hspace{5pt}}c@{\hspace{5pt}}c@{\hspace{5pt}}c@{\hspace{5pt}}c@{\hspace{5pt}}c}
\toprule
Method & \multicolumn{3}{c}{Pinwheel} & \multicolumn{3}{c}{Swiss Roll} & \multicolumn{3}{c}{Square}\\ 
\cmidrule(lr){2-4} \cmidrule(lr){5-7} \cmidrule(lr){8-10}  
& ELBO & MMD & Latent NLL & ELBO & MMD & Latent NLL & ELBO & MMD & Latent NLL \\ 
\midrule
AEVB & $-12.13 \pm 0.41$ & $0.77 \pm 0.04$ & $1.68 \pm 0.31$ & $-14.80 \pm 0.23$ & $0.78 \pm 0.17$ & $5.65 \pm 1.58$ & $-7.85 \pm 0.29$ & $1.10 \pm 0.66$ & $2.78 \pm 0.61$\\ 
AEVB-IAF & $-4.19 \pm 0.05$ & ${0.77 \pm 0.00}$ & $1.64 \pm 0.73$ & $-5.10 \pm 0.30$ & $0.61 \pm 0.15$ & $4.43 \pm 1.09$ & $-3.97 \pm 0.22$ & $0.75 \pm 0.12$ & $1.68 \pm 0.27$\\ 
AAEB & $N/A$ & $0.68 \pm 0.02$ & $1.54 \pm 0.19$ & $N/A$ & $\mathbf{0.52 \pm 0.03}$ & ${3.34 \pm 0.16}$ & $N/A$ & $0.80 \pm 0.02$ & $2.46 \pm 0.46$ \\ 
H-AEVB & $-7.03 \pm 3.13$ & ${0.74 \pm 0.02}$ & $2.25 \pm 3.02$ & $-7.21 \pm 4.62$ & ${0.70 \pm 0.22}$ & $4.04 \pm 4.62$ & $-5.71 \pm 3.05$ & ${0.76 \pm 0.21}$ & ${2.22 \pm 2.03}$\\
PIS & $-7.83 \pm 0.64$ & $0.75 \pm 0.14$ & $6.50 \pm 1.11$ & $-9.83 \pm 0.61$ & $0.61 \pm 0.03$ & $\mathbf{2.40 \pm 1.01}$ & $-7.06 \pm 0.06$ & $0.77 \pm 0.04$ & $3.67 \pm 0.08$\\
DDVI & $\mathbf{-3.88 \pm 0.96}$ & $\mathbf{0.67 \pm 0.04}$ & $\mathbf{1.27 \pm 0.21}$ & $\mathbf{-5.03 \pm 0.58}$ & $0.62 \pm 0.33$ & $3.86 \pm 0.17$ & $\mathbf{-3.79 \pm 0.14}$ & $\mathbf{0.66 \pm 0.07}$ & $\mathbf{1.56 \pm 0.09}$\\ 
\bottomrule
\end{tabular}}
\caption{Unsupervised learning on MNIST. We report ELBO, MMD between generated images and test images, and latent negative log-likelihood (Latent NLL) with pinwheel, swiss roll, and square priors.}
% \vspace{-3pt}
\label{tab:unsupervised}
\end{table*}

\section{Experiments}~\label{sec:exp}

% We report our experimental results in this section and analyzed them correspondingly. To assess whether our method of augmenting variational inference with diffusion encoders can fit an expressive posterior and embed meaningful semantics improving dimensionality reduction and visualization or not, 
We compare DDVI with Auto-Encoding Variational Bayes (AEVB) \citep{kingma2013auto}, AEVB with inverse autoregressive flow posteriors (AEVB-IAF) \citep{kingma2016improved}, Adversarial Auto-Encoding Bayes (AAEB) \citep{makhzani2015adversarial}, and Path Integral Sampler (PIS) \cite{zhang2021path} on MNIST \citep{lecun1998mnist} and CIFAR-10 \citep{krizhevsky2009cifar} in unsupervised and semi-supervised learning settings, and also on the Thousand Genomes dataset \citep{siva20081000}. We also compare with Hierachical Auto-Encoding Variational Bayes (H-AEVB) \citep{ranganath2016hierarchical,vahdat2020nvae} in unsupervised setting. We discuss the computational costs of all methods in Appendix D. The priors, model architecture, and training details can also be founded in Appendix H. All results below are reported with 95\% confidence interval using 3 different seeds.
% We also add PCA \citep{wold1987principal}, t-SNE \citep{van2008visualizing}, and UMAP \citep{mcinnes2018umap} for the comparison of clustering and visualization for genotype analysis.

% \yw{probably: datasets, baselines, configurations}
\subsection{Unsupervised learning}
% Unsupervised learning, fundamental for deep latent variable modeling, enables the identification of complex patterns in data without explicit labels, especially since obtaining labeled data can be costly or impractical. Deep latent variable models extract significant latent representations from unlabeled datasets, emphasizing underlying data variations. For an effective unsupervised framework, a nuanced approximate posterior captures broad data distributions and intricate relationships, while a well-aligned prior regularizes the latent space for smooth, generalizable representations. Together, they ensure the framework captures data nuances and provides meaningful latent insights.

% under unsupervised learning settings in terms of classification accuracy and latent negative log-likelihood (NLL). 
We start with synthetic experiments that are aimed at benchmarking the expressivity of diffusion-based posteriors and their ability to improve fitting $p$, a distribution with a complex structured prior, like one might find in probabilistic programming, scientific analysis, or other applications.
We fit a model $p_\theta(\bx,\bz)$ on the MNIST and CIFAR-10 datasets with three priors $p(\bz)$: pinwheel, swiss roll, and square and report our results in Table \ref{tab:unsupervised} and \ref{tab:unsupervised_cifar}. 
The model distribution $p_\theta$ is instantiated by a deep Gaussian latent variable model (DGLVM) with multi-layer perceptrons (MLPs) on MNIST and convolutional neural networks (CNNs) on CIFAR-10 (see the details of model architecture in Appendix G).

Our first set of metrics (ELBO and MMD) seeks to evaluate the learned generative model $p_\theta$ is good. In the ELBO calculation, we average the reconstruction loss across image pixels.
We use MMD to measure sample quality: we generate images with the trained model and calculate MMD between the generated images and test images using a mixture of Gaussian kernel with sigma equal to $[2, 5, 10, 20, 40, 80]$. We only report MMD for MNIST, since CIFAR-10 generated samples are very low-quality for all methods because the latent dimension is 2.

Our last metric seeks to directly evaluate the expressivity of the posterior.
We measure latent negative log-likelihood (Latent NLL) by fitting a kernel density estimator (KDE) on the latents produced by the model with test data as input and compute the log-likelihood of the latents sampled from the prior under the fitted KDE.

From Table \ref{tab:unsupervised} and Table \ref{tab:unsupervised_cifar} in Appendix, we see our method DDVI achieve best ELBO in all but one scenario, in which it still performs competitively. We also see strong results in Latent NLL and k-nearest neighbors classification accuracy of the latents (Acc). in many scenarios, except for swiss roll where AAEB does well. 
We present visualizations on MNIST using the baselines and our method in Figure \ref{fig:unsupervised}.

\begin{figure*}[t]
    \centering
    % \includegraphics[width=1\textwidth]{imgs/semi-supervised_fig_cropped2.pdf}

    \includegraphics[width=0.95\textwidth]{imgs/semi-supervised_no_font_cropped.pdf}
    %\includegraphics[scale=1.5]{spiral}
    % \vspace{-2em}
    \caption{Semi-supervised visualization on MNIST with 1,000 labels using three different priors (pinwheel, swiss roll, and square). Each a indicates one class.}
    % \vspace{-10pt}
    \label{fig:semi-supervised}
\end{figure*}


\begin{table*}[ht]
\centering
\adjustbox{max width=1\textwidth}{
\centering
\setlength\tabcolsep{4pt} % control the horizontal spacing between columns
\begin{tabular}{c@{\hspace{5pt}}c@{\hspace{5pt}}c@{\hspace{5pt}}c@{\hspace{5pt}}c@{\hspace{5pt}}c@{\hspace{5pt}}c@{\hspace{5pt}}c@{\hspace{5pt}}c@{\hspace{5pt}}c}
\toprule
Method & \multicolumn{3}{c}{Pinwheel} & \multicolumn{3}{c}{Swiss Roll} & \multicolumn{3}{c}{Square}\\ 
\cmidrule(lr){2-4} \cmidrule(lr){5-7} \cmidrule(lr){8-10}  
& ELBO & Acc & Latent NLL & ELBO & Acc & Latent NLL & ELBO & Acc & Latent NLL \\ 
\midrule
AEVB & $-11.15 \pm 0.53$ & $0.93 \pm 0.01$ & $1.36 \pm 0.03$ & $-15.29 \pm 1.33$ & $0.68 \pm 0.01$ & $4.60 \pm 0.23$ & $-10.26 \pm 0.25$ & $0.86 \pm 0.01$ & $1.68 \pm 0.02$\\ 
AEVB-IAF & $-2.10 \pm 0.26$ & $\mathbf{0.95 \pm 0.00}$ & $\mathbf{1.06 \pm 0.03}$ & $-5.38 \pm 1.78$ & $0.90 \pm 0.02$ & $2.75 \pm 0.14$ & $-2.67 \pm 0.83$ & $\mathbf{0.91 \pm 0.01}$ & $\mathbf{0.90 \pm 0.02}$\\ 
AAEB & $N/A$ & $0.89 \pm 0.01$ & $1.55 \pm 0.01$ & $N/A$ & $0.88 \pm 0.01$ & $3.07 \pm 0.05$ & $N/A$ & $1.94 \pm 0.38$ & $0.76 \pm 0.13$ \\ 
DDVI & $\mathbf{-0.24 \pm 0.13}$ & $\mathbf{0.95 \pm 0.00}$ & $\mathbf{1.06 \pm 0.01}$ & $\mathbf{-2.89 \pm 0.33}$ & $\mathbf{0.92 \pm 0.01}$ & $\mathbf{2.09 \pm 0.00}$ & $\mathbf{0.02 \pm 0.09}$ & $0.90 \pm 0.01$ & $1.49 \pm 0.03$\\ 
\bottomrule
\end{tabular}}
\caption{Semi-supervised learning on MNIST (1,000 labels). We report ELBO, accuracy using KNN (K=20) classifier (Acc), and latent negative log-likelihood (Latent NLL) with pinwheel, swiss roll, and square priors.}
% \vspace{-2pt}
\label{tab:semi-supervised}
\end{table*}

\subsection{Semi-supervised Learning}

We also evaluate the performance of our method and the baselines under semi-supervised learning setting where some labels are observed (1,000 for MNIST and 10,000 for CIFAR-10) and the partitions of the priors are known.

For this setting, we evaluate ELBO, Latent NLL, and Acc. We choose classification accuracy since classification is a common downstream task for semi-supervised learning. We use the same set of priors and baselines. Details on how we partition each prior into $p_\btheta(\bz|\bx, l)$ can be founded in Appendix F. The partitions defined for our priors are local parts of the priors. We note that unlike unsupervised learning, we use the simplified sleep term in our objective for this setting (see Appendix B for details), since $q_\phi$ already gets extra information from $l$ here.

The results are shown in Table \ref{tab:semi-supervised} and Table \ref{tab:semi-supervised_cifar} in Appendix. DDVI mostly outperforms the baselines across different priors and metrics, especially on CIFAR-10 where DDVI is best across the board. 
For MNIST, DDVI always achieves the best ELBO, and it also performs competitively with other baselines in classification accuracy.
We also show the visualizations of the latents in Figure \ref{fig:semi-supervised} where DDVI matches the prior almost perfectly.
\begin{table}[ht]
\centering
\resizebox{\linewidth}{!}{
\begin{tabular}{cccc}
\toprule
Method & Cluster Purity & Cluster Completeness & NMI\\
\midrule
AEVB & $0.28 \pm 0.02$ & $\mathbf{0.78 \pm 0.16}$ & $0.59 \pm 0.08$\\
AEVB-IAF & $0.29 \pm 0.04$ & $\mathbf{0.73 \pm 0.06}$ & $0.55 \pm 0.06$\\
AAEB & $0.37 \pm 0.06$ & $\mathbf{0.76 \pm 0.11}$ & ${0.63 \pm 0.02}$\\
DDVI & $\mathbf{0.45 \pm 0.03}$ & $\mathbf{0.75 \pm 0.05}$ & $\mathbf{0.66 \pm 0.04}$\\
\bottomrule
\end{tabular}}
\caption{Quantitative genotype clustering results.}
\label{tab:genotype_clustering}
\end{table}

\subsection{Clustering and~Visualization for Genotype Analysis}

\begin{figure*}[ht]
    \centering
    % \includegraphics[width=\textwidth]{imgs/genotype.pdf} % Adjust the width as needed
    % \vspace{-1em} % Adjust the vertical space as needed

    \includegraphics[width=\textwidth]{imgs/genotype_no_font.pdf}
    \caption{Visualization of genotype clusters. A color represents one ethnicity.}
    \label{fig:genotype}
\end{figure*}

% Clustering helps in revealing the inherent groupings in the data, which can provide insights into data patterns and relationships that might not be immediately evident. By grouping similar data points together, clustering can simplify complex datasets, making them more tractable for analysis or visualization. Moreover, clustering in the latent space facilitates a more organized exploration of data, allowing for insights that might otherwise remain obscured in the high-dimensional original space. Here we unveil our model's capability on whether it can efficiently uncover and delineate the inherent groups or subspaces within the data, leading to a more interpretable and meaningful latent representation. 
In this section, we report results on an real-world task in genome analysis. Visualizing genotype data reveals patterns in the latent ancestry of individuals. 
We instantiate DDVI with a deep Gaussian latent variable model (DGLVM) and compare it against with the three strong clustering baselines using the 1000 Genomes dataset. 
We also report visualizations from three dimensionality reduction algorithms: PCA, TSNE, and UMAP.
For each clustering algorithm, we seek to discover up to 20 clusters. We report quantitative results in terms of cluster purity, cluster completeness, and normalized mutual information (NMI). There is an inherent trade-off between cluster purity completeness. The overall clustering performance can be captured with NMI. 

In Table \ref{tab:genotype_clustering}, we see that DDVI attains the best performance on cluster purity and NMI. For cluster completeness, VAE and AAE have better means but much larger confidence interval.
Furthermore, we visualize our genotype clustering results in latent space, shown in Figure \ref{fig:genotype}, and also report results from classical dimensionality reduction and visualization methods that do not perform clustering (PCA \citep{wold1987principal}, t-SNE \citep{van2008visualizing}, and UMAP \citep{mcinnes2018umap}). The legend of Figure \ref{fig:genotype} can be founded at Figure \ref{fig:legend} in Appendix.

% our DD-VAE can map the latents into fine-grained clusters. In many scenarios, practitioners would want to impose certain structure on the visualization that reflects their prior knowledge on the data. In our experiment, we model such a prior knowledge as a mixture of Gaussians. We can see that DD-VAE outperforms the other methods, including both classic methods and latent variable models.

% To demonstrate our method on a practical task, we experiment on genotype visualization. The goal here is to investigate the population structure in the data. Good genotype visualization methods should visualize known population structures and also reveal previously unknown cryptic structures. In \Cref{fig:viz_genotype}, we compare our method with state-of-the-art visualization methods on 1kgenome visualization. [For this figure, I think we can use our method with the soft learnable prior]

% \input{figs/vizmap}

% In addition, unlike unsupervised visualization methods like t-SNE and UMAP, autoencoder-based methods like ours can perform semi-supervised visualization taking advantage of label information of the data if available. To showcase this capability, we perform semi-supervised visualization using rough map locations of the genotype in \Cref{fig:vizmap}

% \subsection{Visualization}
% \label{sec:exp:viz}

% % We are interested in two visualization scenarios: 
% % \begin{enumerate}
% %     \item Generic visualizations: we want a semantically meaningful visualization; similar datapoints are close to each other and dissimilar are datapoints far apart. 
% %     \item Shape-based visualizations: in addition to producing semantically meaningful visualization, we also want to impose a distribution (or shape) on the latents.
% % \end{enumerate}

% \input{figs/viz}
% \input{figs/viz_toy}

% \Cref{fig:viz} compares our method with other autoencoder-based methods designed to work with non-Gaussian priors on unsupervised and semi-supervised visualization tasks on MNIST. For unsupervised visualization, while other autoencoder-based methods fail to produce prior-matching latent when the priors are sophisticated (like a swiss roll), our method successfully produces latent that match the prior while keeping latent from the same class close to each other. For semi-supervised visualization, we assume that the prior can be partitioned into subpriors, each correspond to different label, and some datapoints we want to visualize are labeled. In this case, our method seems to be able to leverage the label information better in more sophisticated priors. 

% We further investigate visualization of synthetic data that shares the topological/geometrical structure with the prior in \Cref{fig:viz_toy}. We created 3D Swiss roll and square data and visualized it with Swiss roll and square priors respectively. Even when the data and the prior share topological/geometrical structure, only our method produces latent that preserve the structure.

% \Top{Do we also want to include soft priors, i.e., prior with learnable parameters, that I have in this paper? For example, if we have MoG as a prior, we can learn the mean and variance of each Gaussian. Or can we think of any priors whose parameters we want to learn? I think we'll need soft priors for genotype visualization, and it makes sense to introduce it here}

% \subsection{Representation Learning}

% \subsubsection{Latent Quality}
% Good representations from autoencoder-based methods must:
% \begin{enumerate}
%     \item fit the prior well
%     \item provide useful latents for downstream tasks such as classification
% \end{enumerate}

% Thus, we compare our method with other autoencoder-based methods on the divergence between the aggregated posterior and the prior, and the classification accuracies for logistic regression on MNIST and CIFAR10 in \Cref{tab:latent_quality}  

% \Top{Consider including figure of generated samples}

% \input{figs/dis}

% \subsubsection{Disentanglement}

% Generative models have been able to disentangle the class label information from many other latent factors of variation in a principled way. Here, we would like to investigate if our method has such a characteristic. We trained our models on MNIST (FMNIST, 3Dshapes, etc.) in which the latent representation is forced into a N-d Gaussian. We show the disentanglement results qualitatively in \Cref{fig:dis} \todo{Yingheng: which paradigm should we consider here: supervised, semi-supervised, unsupervised?}. We also quantitatively measure the disentanglement performance by calculating DCI scores over latent representations, as shown in \Cref{tab:dis}.

% \subsubsection{Stability/Robustness}

% We hypothesize that our method should be much less sensitive to hyperparameters compared to adversarial-based methods like AAE. In \Cref{tab:hyperparameter_sensitivity}, we show how using the best hyperparameters for MNIST results in suboptimal performance from AAE, while our method's performance does not drop as much.

% \subsection{Clustering}

% Another feature of AAE is clustering. The cluster heads are learnable, and each data point belongs to only one cluster. As mentioned in section ..., we can learn cluster membership $y$ without using adversarial training. In \Cref{tab:clustering}, we show that our method can achieve comparable unsupervised and semi-supervised clustering performance to AAE. 

% \Top{Consider including figure of samples in each cluster}

% \subsection{Genotype Visualization}

% \input{figs/viz_genotype}

% To demonstrate our method on a practical task, we experiment on genotype visualization. The goal here is to investigate the population structure in the data. Good genotype visualization methods should visualize known population structures and also reveal previously unknown cryptic structures. In \Cref{fig:viz_genotype}, we compare our method with state-of-the-art visualization methods on 1kgenome visualization. [For this figure, I think we can use our method with the soft learnable prior]

% \input{figs/vizmap}

% In addition, unlike unsupervised visualization methods like t-SNE and UMAP, autoencoder-based methods like ours can perform semi-supervised visualization taking advantage of label information of the data if available. To showcase this capability, we perform semi-supervised visualization using rough map locations of the genotype in \Cref{fig:vizmap}

% \newpage
% \input{tables/unsup}
% % \input{tables/latent_quality_mnist}
% % \input{tables/latent_quality_cifar}
% \input{tables/clustering}
% \input{tables/dis}
% \input{tables/hyperparameter_sensitivity}
\section{Discussion}

% Our method relates to previous work that uses variational posteriors based on normalizing flows \citep{rezende2015variational} or adversarial networks \citep{makhzani2015adversarial}.

\paragraph{Diffusion vs.~Normalizing Flows}

Our approach is most similar to flow-based approximators \citep{rezende2015variational,kingma2016improved}; in fact when $T \to \infty$, our diffusion-based posterior effectively becomes a continuous-time normalizing flow \citep{song2020denoising}. However, classical flow-based methods require invertible architectures for each flow layer: this constrains their expressivity and requires backpropagating through potentially a very deep network.
%
Our approach, on the other hand, trains a model (a continuous-time flow when $T \to \infty$) via a denoising objective (similar to score matching) that does not require invertible architectures and effectively admits an infinite number of layers (with weight sharing). This model is trained not by backpropagating through the ELBO, but rather via an auxiliary diffusion loss term (effectively, a score matching objective). 

Despite training with a modified loss, we observe in Section \ref{sec:exp} that a diffusion model with an expressive denoising architecture yields an improved ELBO relative to regular flows. Also, our modified loss based on the forward KL divergence reduces posterior collapse (i.e., all modes of the prior are covered well), and thus produces better samples.

\paragraph{Diffusion vs.~Other Generative Models}

Variational posteriors based on GANs \citep{makhzani2015adversarial} also admit expressive architectures and require only sample-based access to the prior $p(\bz)$. Our diffusion-based approach admits a more stable loss, and is potentially more expressive, as it effectively supports an infnite number of layers (with shared parameters when $T\to \infty$). Unlike GANs, our models also admit explicit likelihoods and allow us to compute the ELBO for model evaluation.
Our approach is similar to variational MCMC \citep{salimans2015mcmc}; however, we train with a better objective augmented with a diffusion loss, and we adopt improved architectures with shared weights across layers.

\paragraph{Diffusion for Approximate Inference}
Existing diffusion-based approximate inference methods \cite{berner2022optimal, zhang2021path, vargas2023denoising, zhang2023diffusion, richter2023improved, sendera2024diffusion, akhound2024iterated} 
focus on the task of drawing samples from unnormalized distributions $\tilde p(\bz)$ and estimating the partition function $Z = \int_\bz \tilde p(\bz) d\bz$. 
While these methods are applicable in our setting---we set the unnormalized $\tilde p(\bz)$ to $p_\theta(\bx,\bz)$ such that $Z = p_\theta(\bx)$---they also tackle a more challenging problem (drawing samples from energy-based models) in more general classes of models (arbitrary unnormalized distributions).
%
In contrast, we focus on restricted but still important classes of models (VAEs, Bayes networks, etc.), and we solve more challenging sets of tasks (e.g., maximum-likelihood learning) by using properties of $p_\theta$ (the factorization $p_\theta(\bx|\bz)p_\theta(\bz)$ and efficient sampling from $p_\theta$). 

Our algorithms are also simpler. For example, diffusion sampling methods require backpropagating through a sampling process to minimize the reverse $\text{KL} (q_\phi || p_\theta)$, which poses challenges with optimization and credit assignment. Some methods based on Schrodinger bridges require an iterative optimization process generalizing the sinkhorn algorithm or computationally expensive on-policy or off-policy \cite{malkin2022trajectory} trajectory-based optimization. In contrast, DDVI optimizes the forward $\text{KL} (p_\theta || q_\phi)$ using simple gradient-based optimization that directly emulates diffusion-based training.

% \paragraph{Towards Regularized Variational Inference}\label{sec:regularization}

% Lastly, our method can be understood as introducing a new form of regularization. While a standard VAE introduces a prior regularization term, we introduce a new form of regularization via a function $r$. 
% This can be advantageous when specifying or enforcing certain priors is difficult. 

% For example, the $r$ in 
% $\KL(q_\phi(\by,\bz|\bx)||r(\by|\bx,\bz)p(\bz))$ 
% could be set to a {\em fully deterministic} mapping parameterized with a bijective flow from a Gaussian $\by$ to an arbitrary prior $\bz$ (defined by the flow). The $q_\phi(\bz|\by,\bx)$ that minimizes the KL term is equal to the reverse $r(\bz|\by,\bx)$ of this flow. Thus, for suitably designed $r$, we could try to match the shape of an arbitrary prior specified by $r$, a similar type of regularization could be enforced in the sleep term as well.
% Our method can thus be seen as a probabilistic extension of the InteL-VAE framework \citep{miao2021incorporating}.

% \VK{can cut this more if needed}

% Our framework also supports the user defining general distributions as regularizers $r_1(\bz,\by)$ and $r_2(\by|\bz,\bx)$. Here, $r_1$ puts a custom prior on $\bz,\by$, while $r_2$ specifies a mapping between $\by, \bz$ that should be followed by the encoder model. We write $\Lc_\text{reg}(\bx, \theta, \phi, r_1)$ and $\Lc_\text{sleep}(\bx, \phi, r_2)$, where each term is still defined as in \Cref{eqn:objective1} and \Cref{eqn:sleep1b}, except using $r_1, r_2$ and we optimize


% \subsection{Towards Regularized Variational Inference}\label{sec:regularization}

% Lastly, our approach can be understood from a third perspective: that of regularization. As we mentioned earlier, $\Lc$ is the ELBO with an additional regularization term that encourages the approximate posterior $q(\bz|\bx,\by)q(\by|\bx)$ to gradually transform $\by$ into $\bz$ following the guidance of $r$. In that sense, $r$ provides data augmentation to ensure that the mapping from $\by$ to $\bx$ is easier to learn, similarly to how a diffusion process provides additional learning signal for generating $\bx$ from noise.

% Our framework also supports the user defining general distributions as regularizers $r_1(\bz,\by)$ and $r_2(\by|\bz,\bx)$. Here, $r_1$ puts a custom prior on $\bz,\by$, while $r_2$ specifies a mapping between $\by, \bz$ that should be followed by the encoder model. We write $\Lc_\text{reg}(\bx, \theta, \phi, r_1)$ and $\Lc_\text{sleep}(\bx, \phi, r_2)$, where each term is still defined as in \Cref{eqn:objective1} and \Cref{eqn:sleep1b}, except using $r_1, r_2$ and we optimize
% % The resulting objective becomes
% \begin{align}
%     \Lc(\bx,\theta,\phi,r_1,r_2) = \Lc_\text{rec}(\bx, \theta, \phi) + \beta \cdot \Lc_\text{reg}(\bx, \theta, \phi, r_1) + \gamma \cdot \Lc_\text{sleep}(\bx, \phi, r_2),
% \end{align}
% where $\beta, \gamma > 0$ (in particular, $\beta$ mirrors the $\beta$-VAE model \citep{higgins2016beta}).

% % When using diffusion models, we define our final loss to be:
% % $$
% % \mathcal{L}_d(x, \theta, \phi)
% % = \mathcal{L}(x, \theta, \phi)_\text{wake} + \mathcal{L}(x, \phi)_\text{diff}
% % $$
% % The entire objective is optimized end-to-end

% Examples of $r_2$ include Gaussian diffusion, discrete noising processes, as well as noise conditioned on $\bx$. 
% Note that $r_2$ could even be {\em fully deterministic} when parameterized with a bijective flow. With a strong regularizer, we could set $q(\bz|\by,\bx)$ to equal the reverse $r_2(\bz|\by,\bx)$ of this flow. Thus, for suitably designed $r_2$, we could perfectly match the shape of any mapping. This mirrors the InteL-VAE framework \citep{miao2021incorporating}, with support for flexible and probabilistic regularizers.

% Most importantly, our method can be understood as introducing a new form of regularization. While a standard VAE introduces a prior regularization term, we introduce a new form of regularization via a function $r_2$. 
% % In the extreme case (including in several of our experiments), we can have no prior regularization, and regularize only via a function. 
% This can be advantageous when specifying or enforcing certain priors is difficult. 
% % In our framework, we could define certain functions that would allow enforcing certain prior shapes exactly.



% % \begin{itemize}
% %     \item in the most general case, we may consider two arbitrary regularizers $r_1$, $r_2$
% %     \item We may view this as encoding prior objective not only via a prior, but via a function. In particular, we allow for custom noise processes
% %     \item For example, one choice recovers the previous process. Other choices include generalization to discrete data. Lastly, we could also add conditional noise to the process. These are all supported in our framework.
% % \end{itemize}




\section{Related Work}

% \paragraph{Visualization and Dimensionality Reduction}

% Dimensionality reduction is a cornerstone of machine learning, serving as a vital tool for visualization and data pre-processing. Dimensionality reduction algorithms can be generally bifurcated into two categories. The first emphasizes the preservation of pairwise distance structures among all data samples, with examples being PCA \citep{wold1987principal}, MDS \citep{carroll1998multidimensional}, and Sammon mapping \citep{henderson1997sammon}. In contrast, the second category prioritizes the preservation of local distances over global ones, with t-SNE \citep{van2008visualizing}, Isomap \citep{balasubramanian2002isomap}, LargeVis \citep{tang2016visualizing}, Laplacian eigenmaps \citep{belkin2003laplacian}, diffusion maps \citep{coifman2006diffusion}, and UMAP \citep{mcinnes2018umap} being notable instances.

% In biology, visualizing multi-scale structures is crucial for hypothesizing with insight. As a response, several dimensionality reduction methods tailored for biological data have emerged. Notable among these are PHATE \citep{moon2017phate}, TriMap \citep{amid2019trimap}, and PaCMAP \citep{wang2021understanding}. These methods adeptly balance both global and local structural optimization, paving the way for a more nuanced exploration of the data and fostering the generation of meaningful hypotheses.

% \begin{itemize}
%     \item Identify specific approaches to visualization and dimensionality reduction and cite the most important papers for each approach.
%     \item Cite some algorithms that are specialized to the visualization of biological data.
% \end{itemize}


% \paragraph{Latent Variable Models}
% VAEs \citep{kingma2013auto} enhance the auto-encoder framework by using variational inference algorithms. Many extension of VAEs center on equipping VAEs with adaptable posteriors. InfoVAE \citep{zhao2017infovae} uses mutual information regularization to address a significant VAE drawback, enhancing the variational posterior's quality.  \citet{maaloe2016auxiliary} add auxiliary latent variables for variational distribution flexibility. DVAE \citep{im2017denoising} introduces a unique distribution type, where input noise marginalization over a corruption distribution aids in achieving a flexible approximate distribution class. AAE \citep{makhzani2015adversarial} aligns the latent aggregated posterior distribution with prior distribution using adversarial loss. Flow-based VAEs \cite{rezende2015variational, kingma2016improved} incorporates normalizing flows to build more expressive aggregate posteriors.

% There are a number of recent works merging VAE and diffusion models. 
\paragraph{Latent Diffusion}
\citet{vahdat2021score,wehenkel2021diffusion,rombach2022high} perform diffusion in the latent space of a VAE.
% Note that these latent diffusion models are \textbf{different from our work}. 
Their goal is high sample quality, and they introduce into $p$ hierarchical latents with simple Gaussian priors. 

{Our goal is different}: we seek to fit a $p$ with structured latents (e.g., in probabilistic programming or in science applications, users introduce prior knowledge via hand-crafted $p$), and we improve variational inference in this structured model by introducing auxiliary latents into $q$.
% These works focus on sample quality and directly fit $q_\phi(\bz|\bx)$ to a complex distribution (e.g., compressed images \citep{rombach2022high}); our work instead seeks to achieve high sample quality in latent space, and obtains complex $\bz$ by diffusion from a simple $\by \sim q_\phi(\by|\bx)$.
Recent work \citep{preechakul2022diffusion, zhang2022unsupervised, wang2023infodiffusion} has also melded auto-encoders with diffusion models, focusing on semantically meaningful low-dimensional latents in a diffuser $p$. 
\citet{cohen2022diffusion} crafts a diffusion bridge linking a continuously coded vector to a non-informative prior distribution.
% \vspace{-10pt}
\paragraph{Diffusion for Approximate Inference}

Diffusion sampling \cite{berner2022optimal, zhang2021path, vargas2023denoising, zhang2023diffusion, richter2023improved, sendera2024diffusion, akhound2024iterated} mainly focuses on 
the task of drawing samples from unnormalized distributions and estimating the partition function. 
These works draw connections between diffusion (learning the denoising process) and stochastic control (learning the Föllmer drift).
% leading to several approaches, e.g., path integral sampler (PIS) \cite{zhang2021path}, denoising diffusion sampler (DDS) \cite{vargas2023denoising}, and time-reversed diffusion sampler (DIS) \cite{berner2022optimal}, which have been unified by \citet{richter2023improved}. 
Some other works \cite{zhang2023diffusion, akhound2024iterated, sendera2024diffusion} use continuous generative flow networks (GFlowNets) – deep reinforcement learning algorithms adapted to variational inference that offers stable off-policy training and thus flexible exploration. 
% \citet{sendera2024diffusion} benchmarked these previous diffusion-structured amortized inference methods and studied how to improve credit assignment in diffusion samplers, which refers to the propagation of learning signals from the target density to the parameters of earlier sampling steps. 
While this work is applicable to our setting, it does not rely on the structure of $p(\bx,\bz)$ available to us, namely tractable sampling in $p$.

% The $\triangle$VAE [CITE] employs Diffusion Variational Autoencoders, leveraging Brownian motion transition kernels to accommodate an arbitrary manifold in the latent space, reflecting data's topological attributes.

% VAEs [CITE] enhance the auto-encoder framework by using variational inference algorithms, creating a generative model with meaningful and smooth latent spaces. They've been mainly employed for disentangling latent factors, as seen in the $\beta$-VAE [CITE] and factor-VAE [CITE]. The $\triangle$VAE [CITE] employs Diffusion Variational Autoencoders, leveraging Brownian motion transition kernels to accommodate an arbitrary manifold in the latent space, reflecting data's topological attributes. Other works have centered on equipping VAEs with adaptable posteriors. InfoVAE [CITE] uses mutual information regularization to address a significant VAE drawback, enhancing the variational posterior's quality. Salimans et al. (2015) [CITE] incorporate MCMC steps into variational inference, ensuring the variational distribution progressively aligns with the target distribution. Maaloe et al. (2016) [CITE] add auxiliary latent variables for greater variational distribution flexibility. AAE [CITE] merges traditional reconstruction with adversarial training to achieve an implicit generative model that aligns the latent aggregated posterior distribution with any given prior distribution. DVAE [CITE] introduces a unique distribution type, where input noise marginalization over a corruption distribution aids in achieving a flexible approximate distribution class.

% Vahdat et al. (2021) [CITE] and Wenhenkel et al. (2021) [CITE] both recommend using diffusion [CITE] to model VAE prior distributions, enabling enhanced generative modeling. Recent research [CITE] has also melded auto-encoders with diffusion models, focusing on semantically meaningful latents and supporting stable, high-quality training across diverse datasets. Another prevalent strategy involves the normalizing flows framework [CITE]. Chen et al. (2017) [CITE] suggest learning the prior to bridge the gap between approximate and true posteriors, using an autoregressive flow for the former and an inverse autoregressive flow for the latter [CITE]. Tomczak [CITE] introduces the VampPrior, conceptualized from the aggregated posterior [CITE], as a mixture of learned pseudo-inputs. Cohen [CITE] crafts a diffusion bridge linking a continuous coded vector to a non-informative prior distribution, where latent discrete states emerge as random functions of these continuous vectors.

% \begin{itemize}
%     \item Cite recent papers that mix diffusion models and VAEs.
%     \item Cite and compare against papers that do diffusion in latent space (latent diffusion, score-based generative models in latent space, diffusion bridges for VQVAE, etc.)
%     \item Cite general research in generative models, normalizing flows, approximate posteriors
% \end{itemize}

% \section{Discussion and Conclusion}

% \paragraph{Limitations of Diffusion Models as Variational Posteriors} D\paragraph{Limitations of Diffusion Models as Variational Posteriors} iscuss the fact that diffusion models can approximate complex priors, but their learning objective is not the ELBO; hence it may not be as good at other tasks like density estimation. Discuss which tasks it's better on and why.

% \paragraph{Latent Variable Models for Dimensionality Reduction}
% Discuss the pros and cons of using an LVM-based approach relative to something like UMAP or TSNE.

\paragraph{Dimensionality Reduction}
Latent variable models in general are an attractive alternative to visualization methods like PCA, Sparse PCA, NMF, UMAP, and t-SNE \citep{wold1987principal,kuleshov2013fast,NIPS2000_f9d11525,mcinnes2018umap,van2008visualizing}. Domain-specific knowledge can be injected through the prior, and deep neural networks can be utilized to achieve a more expressive mapping from the data space to the latent space. Nevertheless, downsides of LVMs are that they are more computationally expensive and require careful hyperparameter tuning. %Training deep latent variable models requires the use of GPU and takes more than just a few minutes, as opposed to other visualization methods, to complete. Deploying latent variable models on different data and/or prior may also call for different hyperparameters. 

\section{Conclusion}

While this paper focuses on applications of DDVI to dimensionality reduction and visualization, there exist other tasks for the algorithm, e.g., density estimation or sample quality. 
Accurate variational inference has the potential to improve downstream applications of generative modeling, e.g., decision making \citep{DBLP:conf/icml/NguyenG22,deshpande2023calibrated}, meta-learning \citep{rastogi2023semi}, or causal effect estimation \citep{deshpande2022deep}.

Since our learning objective differs from the ELBO (it adds a regularizer), we anticipate gains on models whose training benefits from regularization, but perhaps not on all models. Also, attaining competitive likelihood estimation requires architecture improvements that are orthogonal to this paper. However, our ability to generate diverse samples and achieve class separation in latent space hints at the method's potential on these tasks. 
% Similarly, high class separation in latent space suggests the possibility to improve semi-supervised learning accuracy.
% References
% \bibliographystyle{aaai2025}
\bibliography{reference}


\newpage

\onecolumn

\title{Appendix}
\maketitle

\appendix
\section{Pseudocode}
Here we provide a pseudocode to illustrate the training process of DDVI. 
% We release anonymized code repo at \url{https://anonymous.4open.science/r/vi-with-diffusion-D70F}.
\begin{algorithm*}[h]
\caption{DDVI Pseudocode}
\begin{algorithmic}[1]
\State (Optional) Pre-train $p_\theta(\bx|\bz)$ and $q_\phi(\by|\bx)$ with DDVI but with unconditional diffusion model $q_\phi(\bz|\by)$
\For{epoch $= 1, \ldots, n$}
    \For{$\bx_1, \ldots, \bx_k \sim p_\mathcal{D}(\bx)$}
        \State $\by_i \sim q_\phi(\by|\bx_i)$ and $\bz_i \sim q(\bz|\by_i, \bx_i)$ for $i = 1, \ldots, k$
        \State Optimize $\theta, \phi$ with respect to a Monte Carlo estimate of $\mathbb{E}_{q_\phi(\by,\bz|\bx)} \left[ \log p_\theta(\bx|\bz) \right] - \mathrm{KL}\left(q_\phi(\by,\bz|\bx) \| p_\theta(\by,\bz) \right)$ for each $\bx_i$ \Comment{Standard ELBO training part}
        \For{iteration $= 1, \ldots, m$} \Comment{Do sleep for $m$ iterations}
            \State $\bz_{1}, \ldots, \bz_{k} \sim p(\bz)$ \Comment{Batch-sample latents from prior}
            \State $\hat{\bx}_i \sim p(\bx|\bz_i)$ for $i = 1, \ldots, k$ \Comment{Construct fantasy inputs}
            \State $\by_i \sim r(\by|\bz_i)$ for $i = 1, \ldots, k$ \Comment{Construct fantasy inputs}
            \State Optimize $\phi$ using the standard diffusion noise prediction loss on $q_\phi(\bz|\by_i, \hat{\bx}_i)$
        \EndFor
    \EndFor
\EndFor
\end{algorithmic}
\label{alg:ddvi}
\end{algorithm*}


\section{Simplifying Wake-Sleep}\label{app:sim}
In wake-sleep, sampling $\bx$ from $p_\theta$ to obtain gradients for the sleep term introduces computational overhead. To address this issue, we propose {\em wake-sleep in latent space}, an algorithm that optimizes an approximation $\hat \Lc(\bx,\theta,\phi)$ of $\Lc$:
\begin{align}
    \hat \Lc(\bx, \theta, \phi)
    % & = \Exp_{q_\phi(y|x)} [ \log p_\theta(x|y)  -\KL(q_\phi(y|x)|p_\theta(y)) - \KL(p_\theta(z|y)|q_\phi(z|y))] \label{eqn:elbo_z1} \\
     & =
    \underbrace{\Exp_{q_\phi(\by,\bz|\bx)} [\log p_\theta(\bx|\bz)]}_\text{wake / reconstr. term $\Lc_\text{rec}(\bx,\theta,\phi)$} \underbrace{-\KL(q_\phi(\by,\bz|\bx)||p_\theta(\by,\bz))}_\text{prior regularization term $\mathcal{L}_\text{reg}(\bx,\theta,\phi)$}
    \underbrace{-\KL(p_\theta(\bz)||q_\phi(\bz|\bx))}_\text{latent sleep term $\Lc_\text{sleep}(\bx,\phi)$}.
    \label{eqn:objective2}
\end{align}
We have replaced $\Lc_\text{sleep}(\phi)$ with a latent sleep term $\Lc_\text{sleep}(\bx,\phi)$, in which $\bx$ is given, and we only seek to fit the true reverse noising process $r(\bz|\by)$ independently of $\bx$. We can similarly show that
\begin{align}
\Lc_\text{sleep}(\bx, \phi) 
& = \Exp_{p_\theta(\bz)} [\log {q_\phi(\bz|\bx)}] + \bar H(p_\theta)
 \geq \Exp_{p_\theta(\bz) r(\by|\bz)} [\log ({q_\phi(\by,\bz|\bx)}/r(\by|\bz))] + \bar H(p_\theta) \label{eqn:sleep2} \\
 & = -\Exp_{p_\theta(\bz)} [\KL(r(\by|\bz)||q_\phi(\by|\bz,\bx))] - \KL(p_\theta(\bz)||q(\bz|\bx)), \label{eqn:sleep2_kl}
\end{align}
where $\bar H(p_\theta)$ is an entropy term constant in $\phi$.
Thus, we minimize the forward KL divergence 
% $KL(p_\theta(\bz|\by,\b)|q_\phi(z|y))$ 
by sampling $\bz$, and applying the noising process to get $\by$; the $q_\phi$ is fit to denoise $\bz$ from $\by$ as in Equation \ref{eqn:sleep1b}.


We optimize our bound on $\hat \Lc(x,\theta,\phi)$ end-to-end using  minibatch gradient descent over $\theta, \phi$. 
% This approach jointly optimizes two sets of terms, which we call the conditional wake and sleep terms. The wake term is a regularized reconstruction error: at each input $x$, we encode into a noise variable $y \sim q_\phi(y|x)$, which we transform into $z$ using $q_\phi(z|x)$; we then decode from $z$. 
While the wake term is a reconstruction loss as in wake-sleep, the sleep term generates latent samples $\bz, \by$ from $r(\by|\bz)p_\theta(\bz)$ (by analogy with $p_\theta(\bx|\bz)p_\theta(\bz)$ in normal wake-sleep); the denoiser $q_\phi$ is trained to recover $\bz$ from $\by$.
%
% This learning objective effectively  combines variational inference and wake-sleep. Traditionally, wake-sleep alternates between a wake phase, in which we sample $x$ from the data distribution and minimize the reconstruction error, and a sleep phase, in which $x$ is sampled from the model and we minimize the KL divergence with an approximate prior. 
Thus, we perform wake-sleep {\em in latent space}, 
% and each phase becomes a term: in the wake term, the $y$ variables are sampled from $q(y|x)$, while in the sleep term we $y$ is sampled from $r(y|z)$. Wake sleep in latent space 
which obviates the need for alternating wake and sleep phases, and allows efficient end-to-end training. A limitation of this approximation is that the sleep term does not fit $q_\phi$ to the true $p_\theta(\bz|\bx,\by)$, and as a consequence $\hat L$ is not a tight lower bound on $\log p_\theta(\bx)$. We may think of $\Lc_\text{sleep}(\bx, \phi)$ as a regularizer to the ELBO.

\section{Comparision of Methods}
We provide a comprehensive comparison of different methods in Table \ref{tab:method_comparsion}.
\citet{vahdat2021score,wehenkel2021diffusion,rombach2022high} perform diffusion in the latent space of a VAE to improve the efficiency of image generation.
% Note that these latent diffusion models are \textbf{different from our work}. 
Their goal is high sample quality, and they introduce into $p$ hierarchical latents with simple Gaussian priors. \textbf{Our goal is different}: we seek a method to fit a $p$ with structured latents (e.g., in probabilistic programming or in science applications, users introduce prior knowledge via hand-crafted $p$), and we improve variational inference in this structured model by introducing auxiliary latents into $q$.
% These works focus on sample quality and directly fit $q_\phi(\bz|\bx)$ to a complex distribution (e.g., compressed images \citep{rombach2022high}); our work instead seeks to achieve high sample quality in latent space, and obtains complex $\bz$ by diffusion from a simple $\by \sim q_\phi(\by|\bx)$.

Recent work \citep{preechakul2022diffusion, zhang2022unsupervised, wang2023infodiffusion} has also melded auto-encoders with diffusion models, focusing on semantically meaningful low-dimensional latents in a diffuser $p$. 
\citet{cohen2022diffusion} crafts a diffusion bridge linking a continuous coded vector to a non-informative prior distribution.
\begin{table*}[htbp!]
\adjustbox{max width=\textwidth}{
\centering
\renewcommand{\arraystretch}{1.5}
\setlength\tabcolsep{4pt} % control the horizontal spacing between columns
    \definecolor{OliveGreen}{HTML}{808000}    % Define OliveGreen
    \definecolor{bloodred}{HTML}{B00000}      % Define bloodred
    \definecolor{cautionyellow}{HTML}{EED202} % Define cautionyellow

    % Command to create a colored checkmark
    \newcommand*\colourcheck[1]{%
      \expandafter\newcommand\csname #1check\endcsname{%
        \textcolor{#1}{\ding{51}}% \ding{51} is a checkmark
      }%
    }
    
    % Command to create a colored x-mark
    \newcommand*\colourxmark[1]{%
      \expandafter\newcommand\csname #1xmark\endcsname{%
        \textcolor{#1}{\ding{55}}% \ding{55} is an x-mark
      }%
    }
    
    % Command to create a colored alternative checkmark (if needed)
    \newcommand*\colourcheckodd[1]{%
      \expandafter\newcommand\csname #1checkodd\endcsname{%
        \textcolor{#1}{\ding{53}}% \ding{53} can be another style of checkmark
      }%
    }

    \colourcheck{OliveGreen}       % Defines \OliveGreencheck
    \colourcheck{cautionyellow}    % Defines \cautionyellowcheck
    \colourcheckodd{cautionyellow}% Defines \cautionyellowcheckodd (if needed)
    \colourxmark{bloodred}         % Defines \bloodredxmark
    
    % Create your custom commands using the newly defined commands
    \newcommand{\ourxmark}{\bloodredxmark}%
    \newcommand{\ourcheckmark}{\OliveGreencheck}%

    \newcolumntype{C}[1]{>{\centering\arraybackslash}p{#1}}
\begin{tabular}{C{2cm} C{2.3cm} C{2.6cm} C{2cm} C{1.5cm} C{3cm} C{3.5cm}}
\toprule
Model & Training Objective & Approximating Family & Sample-based Prior & Auxiliary Variable & Tasks & Simplified Graphical Illustration \\
\midrule
AEVB & ELBO & Diagonal Gaussian & \ourxmark & \ourxmark & Density estimation & $\mathbf{x} \rightarrow \mathbf{z} \rightarrow \mathbf{x}$ \\
AEVB-IAF & ELBO & Normalizing flow & \ourxmark & \ourcheckmark & Density estimation / Visualization & $\mathbf{x} \rightarrow \mathbf{z}_0 \rightarrow \mathbf{z}_T \rightarrow \mathbf{x}$ \\
AAEB & Adversarial training & Adversarial generator & \ourcheckmark & \ourxmark & Visualization & $\mathbf{x} \rightarrow \mathbf{z} \rightarrow \mathbf{x}$ \\
H-AEVB-(IAF) & ELBO & Factorial Normal / Normalizing flow & \ourxmark & \ourcheckmark & Density estimation / High-quality sample generation & $\mathbf{x} \rightarrow \mathbf{z}_0 \rightarrow \mathbf{z}_T \rightarrow \mathbf{z}_0 \rightarrow \mathbf{x}$ \\
ADGM & ELBO & Non-Gaussian & \ourxmark & \ourcheckmark & Density estimation & $\mathbf{x} \rightarrow \mathbf{a} \rightarrow \mathbf{z} \rightarrow \mathbf{x}$ \\
LDM & ELBO & Diagonal Gaussian & \ourxmark & \ourcheckmark & High-quality sample generation & $\mathbf{x} \rightarrow \mathbf{z}_0 \rightarrow \mathbf{z}_T \rightarrow \mathbf{z}_0 \rightarrow \mathbf{x}$ \\
LSGM & ELBO \& score matching & Diagonal Gaussian & \ourxmark & \ourcheckmark & High-quality sample generation & $\mathbf{x} \rightarrow \mathbf{z}_0 \rightarrow \mathbf{z}_T \rightarrow \mathbf{z}_0 \rightarrow \mathbf{x}$ \\
DDVI & ELBO \& sleep term & Denoising diffusion & \ourcheckmark & \ourcheckmark & Density estimation / Visualization & $\mathbf{x} \rightarrow \mathbf{z}_T (\mathbf{y}) \rightarrow \mathbf{z}_0 (\mathbf{z}) \rightarrow \mathbf{x}$ \\
\bottomrule
\end{tabular}}
\caption{Comparison of DDVI to other relevant methods. $\mathbf{x}$ represents the original data input to the model. $\mathbf{z}$ denotes the latent (hidden) representation of the input data. $\mathbf{a}$ represents an auxiliary variable introduced in some models (like ADGM) to capture additional aspects of the data distribution or to assist in the model's learning process.}
\label{tab:method_comparsion}
\end{table*}

\section{Computational Cost Analysis}
\label{app:computational_cost}

\begin{table*}[ht]
\adjustbox{max width=\textwidth}{
\centering
\setlength\tabcolsep{4pt} % control the horizontal spacing between columns
\begin{tabular}{ccccccc}
\toprule
Method & \multicolumn{6}{c}{NMI values at different wall-clock training times}\\
\cmidrule{2-7}
& NMI @ 10 min & NMI @ 20 min & NMI @ 30 min & NMI @ 40 min & NMI @ 50 min & NMI @ 60 min \\
\midrule
AEVB             & $0.52$       & $0.52$       & $0.52$       & $0.52$       & $0.52$       & $0.52$       \\
AEVB-IAF         & $0.54$       & $0.52$       & $0.52$       & $0.52$       & $0.52$       & $0.52$       \\
AAEB             & $0.61$       & $0.57$       & $0.57$       & $0.57$       & $0.57$       & $0.57$       \\
DDVI (T=5)   & \textit{warm up} & $0.63$ & $0.63$ & $0.66$ & $0.66$ & $0.66$ \\
DDVI (T=10)  & \textit{warm up} & $\mathbf{0.64}$ & $\mathbf{0.68}$ & $\mathbf{0.70}$ & $\mathbf{0.70}$ & $\mathbf{0.70}$ \\
DDVI (T=20)  & \textit{warm up} & $0.50$ & $0.51$ & $0.56$ & $0.64$ & $0.68$ \\
DDVI (T=50)  & \textit{warm up} & $0.52$ & $0.54$ & $0.51$ & $0.59$ & $0.59$ \\
\bottomrule
\end{tabular}}
\caption{Computational cost trade-off on 1kgenome: NMI vs wall-clock training time}
\label{tab:computational_cost}
\end{table*}

We conduct a computational cost analysis between the baselines and DDVI with various timesteps on the genotype clustering/visualization experiments. Table \ref{tab:computational_cost} shows that DDVI outperforms baselines at all timestamps and continues to improve after the baselines have plateaued.

\section{Connections to Diffusion Samplers}
Diffusion sampling \cite{berner2022optimal, zhang2021path, vargas2023denoising, zhang2023diffusion, richter2023improved, sendera2024diffusion, akhound2024iterated} mainly focuses on 
the task of drawing samples from unnormalized distributions and estimating the partition function. 
These works draw connections between diffusion (learning the denoising process) and stochastic control (learning the Föllmer drift), leading to several approaches, e.g., path integral sampler (PIS) \cite{zhang2021path}, denoising diffusion sampler (DDS) \cite{vargas2023denoising}, and time-reversed diffusion sampler (DIS) \cite{berner2022optimal}, which have been unified by \citet{richter2023improved}. Some other works \cite{zhang2023diffusion, akhound2024iterated} use continuous generative flow networks (GFlowNets) – deep reinforcement learning algorithms adapted to variational inference that offers stable off-policy training and thus flexible exploration. \citet{sendera2024diffusion} benchmarked these previous diffusion-structured amortized inference methods and studied how to improve credit assignment in diffusion samplers, which refers to the propagation of learning signals from the target density to the parameters of earlier sampling steps.  Overall, there are indeed some strong connections between these works and ours:
\begin{itemize}
    \item They also focus on variational methods that directly fit a parametric family of tractable distributions (given by controlled SDEs) to the target density.
    \item They cast the density estimation/sampling problem into an optimization problem over a control objective, which learns control drifts (and diffusion) parameterized by neural networks.
\end{itemize}
But we would like to clarify that there are also some clear differences between them:

\begin{itemize}
    \item The diffusion-structured samplers only focus on density estimation/sampling but ignore the problem of learning a generative model, which is one of the main focuses of our work. We aim to perform more accurate variational inference using an auxiliary variable model augmented by diffusion models to improve generative modeling. In our setting, $p_\theta(z|x)$ is a moving target density, as we jointly learn $\theta$ with $\phi$, as opposed to a static target density that diffusion-structured samplers are designed to solve.
    \item To tackle the challenge of credit assignment – propagating weak learning signals through the sampling trajectory, the techniques proposed in diffusion-structured samplers are mostly based on partial trajectory information, which has higher training costs over on-policy \cite{zhang2021path} or off-policy \cite{malkin2022trajectory} trajectory-based optimization. Instead, we introduce a wake-sleep optimization algorithm and its simplified version to alleviate the weak learning signal issue and optimize the evidence lower bound in a better way.
    \item In Equation \ref{eqn:sleep1b}, we are minimizing the forward KL divergence $\KL (p_\theta || q_\phi)$, where diffusion samplers are minimizing the reverse $\KL (q_\phi || p_\theta)$.
\end{itemize}

We also summarize the connections and differences in the table below.

\begin{table}[ht]
\centering
\begin{tabular}{@{}p{2cm}p{4cm}p{4cm}p{4cm}@{}}
\toprule
                        & \textbf{Diffusion Samplers} & \textbf{GFlowNet-based} & \textbf{DDVI (ours)} \\ 
                        & \cite{berner2022optimal, zhang2021path, vargas2023denoising, akhound2024iterated}              & \cite{zhang2023diffusion, richter2023improved, sendera2024diffusion, malkin2022trajectory}               &             \\ \midrule
\textbf{Tasks}                   & Sampling, density estimation  & Sampling, density estimation  & Learning, Sampling, Dimensionality reduction \\
\textbf{Model Family for $p$}    & Any energy-based              & Any energy-based              & Latent with tractable $p(x|z), p(z)$ \\
\textbf{Model Family for $q$}    & Markov chain                  & Markov chain                  & Markov chain \\
\textbf{Objective}  & $\KL(q||p)$ with regularizer & Trajectory balance objective & $\KL(p||q)$ with ELBO \\
\textbf{Algorithm}  & Gradient descent (with reference process), importance sampling & RL-motivated off-policy optimization (replay buffers, Thompson sampling, etc.) & Gradient descent with wake-sleep \\
\textbf{Compatible Models} & Anything energy-based  & Anything energy-based  & LDA, deep latent-variable models \\
\textbf{Applications}            & Sampling from physics-based models, model selection based on NLL & Sampling from physics-based models, model selection based on NLL & Probabilistic programming, visualization \\ \bottomrule
\end{tabular}
\caption{Comparison of Diffusion-structured Samplers, GFlowNet-based Approaches, and DDVI}
\label{tab:comparison}
\end{table}

\section{Priors}
\label{app:priors}

Below we describe the sampling process for each prior. 

\textbf{Pinwheel.} This distribution was used in \citep{johnson2016svae}. We define the number of clusters to be 10. For semi-supervised learning experiments, this prior is partitioned into 10 partitions, each partition being a cluster.

\textbf{Swiss Roll.} This distribution was used in \cite{marsland2014swissroll}. For semi-supervised learning experiments, this prior is partitioned into 10 partitions. The samples from the prior can actually be characterized by a single scalar representing how far you are long the swiss roll from the center. The paritioning is done by creating 10 equal-length intervals in this 1D space.

\textbf{Square.} This distribution has the shaped of a square going from -1 to 1 in both axes. Each position on the square can be characterized by a single scalar representing how far you are from the top left corner. Sampling is done by sampling the position uniformly and turn the 1D position to 2D latent. We add noise $\sigma = 0.06$ to the prior. For semi-supervised learning experiments, this prior is partitioned into 10 partitions. The partitioning is done by creating 10 equal-length intervals in the 1D position space.

AEVB and AEVB-IAF requires that we can evaluate the prior density. To do this, for all priors, we evaluate the density by fitting a kernel density estimator with mixture of gaussian kernel with bandwidth equal to 0.005, 0.008, 0.01, 0.03, and 0.05. 


\section{Model Architecture}
\label{app:model_architecture}

All methods use the same architecture for encoder $q_\phi(\bz|\bx)$ and decoder $p_\theta(\bx|\bz)$, excluding the extra parts specific to each method which we describe below, for the same dataset. For MNIST, the encoder and decoder are multi-layer perceptron with two hidden layers, each with 1000 hidden units. For CIFAR-10, the encoder is a  4-layer convolutional neural network with (16, 32, 64, 128) channels with a linear layer on top, and the decoder is a 4-layer tranposed convolutional neural network with (64, 32, 16, 3) channels where a linear layer is used to first turn the feature dimension from 2 to 64. 

\textbf{AEVB-IAF} employs 4 IAF transformations on top of the encoder, each is implemented with a 4-layer MADE. The number of hidden units in MADE is 128. The ordering is reversed between every other IAF transformation. 

\textbf{AAEB} has a discriminator, used in adversarial training, which is a multi-layer perceptron with two hidden layers, each with 1000 hidden units.

\textbf{DDVI} has a diffusion model on top of the encoder. The time-conditioned reverse diffusion distribution is implemented with a 5-layer time-conditioned multi-layer perceptron, each with 128 hidden units. A time-conditioned linear layer learns an additional embedding for each timestep and adds it to the output of the linear layer.

\section{Training Details} 
\label{app:training_details}

For training, we update the parameters for each batch of inputs by alternating between the ELBO phase (optimizing $\theta$ and $\phi$ with respect to the ELBO, i.e., the reconstruction term and the prior matching term) and the sleep phase (optimizing $\phi$ with respect to the sleep term). We use Adam optimizer and latent size of 2 for all of our experiments. Each algorithm takes roughly 2 hours on a single Nvidia GeForce RTX 3090 to complete one run of experiment. The training details of each algorithm are detailed below:

\textbf{AEVB.} The batch size is set to 128. The number of epochs is 200 for unsupervised and clustering experiments and 50 for semi-supervised experiments. The learning rate is 0.0001. The loss is BCE for MNIST and CIFAR-10 experiments and MSE for genotype analysis experiments. For semi-supervised MNIST experiments, the kl divergence weight is set to be 0.01, while for semi-supervised CIFAR-10 experiments, the kl divergence weight is set to be 0.01. For other experiments, the KL divergence weight is set with a schedule linear on number of epochs going from 0 to 0.01. We also have a weight of 5 multiplied to the prior density.

\textbf{AEVB-IAF.} The batch size, number of epochs, learning rate, loss, KL divergence weight, and prior density weight are the same as VAE. The context size, i.e., the size of features used to initialize the flow layers for different datat point, is 10. 

\textbf{AAEB.} The batch size is set to 128. The number of epochs is 200 for all experiments. The learning rate is 0.0002. The loss is MSE for all experiments. To stabilize the training, we add noise to the input to the discriminator with sigma 0.3 at the start and lower it by 0.1 for every 50 epochs. The noise equals to 0 at epoch 150.

\textbf{DDVI.} The batch size is set to 128 for most experiments, except for semi-supervied experiments where the batch size is 1024. The number of epochs is 200 for unsupervised and clustering experiments and 30 for semi-supervised experiments. The learning rate is 0.0001. The loss is BCE for MNIST and CIFAR-10 experiments and MSE for genotype analysis experiments. For unsupervised MNIST and CIFAR-10 experiments, the KL divergence weight is set to 0.003. For semi-supervised MNIST experiment, we use KL divergence weight of 0.1. For semi-supervised CIFAR-10 experiment, we use KL divergence weight of 0.5. For clustering experiment, we use KL divergence weight of 0.005. The number of timesteps is 20 for unsupervised and clustering experiments and 100 for semi-supervised experiments.

\section{Genotype Analysis Experiments Details}

\begin{figure*}[t]
    \centering
    \includegraphics[width=1\textwidth]{imgs/1kgenome_legend2_cropped.pdf}
    \vspace{-1em}
    \caption{Legend showing what ethnicity each color corresponds to in the 1000 Genomes dataset}
    \vspace{-10pt}
    \label{fig:legend}
\end{figure*}

Before inputting the data points into any of the visualization methods, we first pre-process it by running a PCA and keep only the first 1000 principal components of the data points. We further divide the features by 30 for all latent variables model methods. 

The legend of the 1000 Genomes Visualization plot can be found at Figure \ref{fig:legend}.

\begin{table*}[ht]
\centering
% \adjustbox{max width=\textwidth}{
\centering
\setlength\tabcolsep{4pt} % control the horizontal spacing between columns
\begin{tabular}{c@{\hspace{5pt}}c@{\hspace{5pt}}c@{\hspace{5pt}}c@{\hspace{5pt}}c@{\hspace{5pt}}c@{\hspace{5pt}}c}
\toprule
Method & \multicolumn{2}{c}{Pinwheel} & \multicolumn{2}{c}{Swiss Roll} & \multicolumn{2}{c}{Square}\\ 
\cmidrule(lr){2-3} \cmidrule(lr){4-5} \cmidrule(lr){6-7}  
& ELBO & Latent NLL & ELBO & Latent NLL & ELBO & Latent NLL \\ 
\midrule
AEVB & $-12.96 \pm 1.81$ & $3.26 \pm 0.60$ & $-12.87 \pm 4.55$ & $6.25 \pm 1.58$ & $-7.91 \pm 0.11$ & $2.91 \pm 0.17$\\ 
AEVB-IAF & $-3.24 \pm 0.16$ & $1.71 \pm 0.84$ & $-4.03 \pm 0.73$ & $5.51 \pm 0.51$ & $\mathbf{-2.10 \pm 0.31}$ & $1.71 \pm 0.77$\\
AAEB & $N/A$ & $1.70 \pm 0.41$ & $N/A$ & $\mathbf{3.18 \pm 0.22}$ & $N/A$ & $1.67 \pm 0.17$\\
H-AEVB & $-4.42 \pm 0.46$ & $\mathbf{1.69 \pm 0.17}$ & $-5.36 \pm 0.77$ & $5.74 \pm 0.55$ & $-2.86 \pm 0.11$ & $1.64 \pm 0.09$\\
PIS & ${-2.92 \pm 1.23}$ & ${3.61 \pm 0.62}$ & $-4.14 \pm 0.49$ & $7.14 \pm 0.14$ & $-4.85 \pm 0.06$ & $3.91 \pm 0.06$\\
DDVI & $\mathbf{-1.38 \pm 0.44}$ & ${1.75 \pm 0.53}$ & $\mathbf{-3.05 \pm 0.65}$ & ${5.66 \pm 2.63}$ & $-2.47 \pm 0.30$ & $\mathbf{1.58 \pm 0.09}$\\

\bottomrule
\end{tabular}
\caption{Unsupervised learning on CIFAR-10. We report ELBO and latent negative log-likelihood (Latent NLL) with pinwheel, swiss roll, and square priors.}
\label{tab:unsupervised_cifar}
\end{table*}
\begin{table*}[ht]
\centering
\adjustbox{max width=\textwidth}{
\centering
\setlength\tabcolsep{4pt} % control the horizontal spacing between columns
\begin{tabular}{c@{\hspace{5pt}}c@{\hspace{5pt}}c@{\hspace{5pt}}c@{\hspace{5pt}}c@{\hspace{5pt}}c@{\hspace{5pt}}c@{\hspace{5pt}}c@{\hspace{5pt}}c@{\hspace{5pt}}c}
\toprule
Method & \multicolumn{3}{c}{Pinwheel} & \multicolumn{3}{c}{Swiss Roll} & \multicolumn{3}{c}{Square}\\ 
\cmidrule(lr){2-4} \cmidrule(lr){5-7} \cmidrule(lr){8-10}  
& ELBO & Acc & Latent NLL & ELBO & Acc & Latent NLL & ELBO & Acc & Latent NLL \\ 
\midrule
AEVB & $-17.14 \pm 1.46$ & $0.30 \pm 0.05$ & $2.32 \pm 0.27$ & $-17.89 \pm 5.21$ & $0.20 \pm 0.07$ & $6.56 \pm 2.25$ & $-13.30 \pm 1.50$ & $0.30 \pm 0.05$ & $1.95 \pm 0.28$\\ 
AEVB-IAF & $-5.70 \pm 0.07$ & $0.47 \pm 0.01$ & $1.62 \pm 0.05$ & $-5.53 \pm 2.82$ & $0.28 \pm 0.08$ & $6.82 \pm 1.90$ & $-4.41 \pm 0.53$ & $0.36 \pm 0.01$ & $1.58 \pm 0.15$\\ 
AAEB & $N/A$ & $0.25 \pm 0.01$ & $1.77 \pm 0.14$ & $N/A$ & $0.23 \pm 0.01$ & $3.38 \pm 0.30$ & $N/A$ & $0.23 \pm 0.04$ & $1.74 \pm 0.15$ \\ 
DDVI & $\mathbf{-1.60 \pm 0.29}$ & $\mathbf{0.49 \pm 0.01}$ & $\mathbf{1.09 \pm 0.05}$ & $\mathbf{-4.13 \pm 1.51}$ & $\mathbf{0.47 \pm 0.09}$ & $\mathbf{2.29 \pm 0.08}$ & $\mathbf{-1.73 \pm 0.64}$ & $\mathbf{0.49 \pm 0.01}$ & $\mathbf{1.48 \pm 0.02}$\\ 
\bottomrule
\end{tabular}}
\caption{Semi-supervised learning on CIFAR-10 (10,000 labels). We report ELBO, accuracy using KNN (K=20) classifier (Acc), and latent negative log-likelihood (Latent NLL) with pinwheel, swiss roll, and square priors.}
\label{tab:semi-supervised_cifar}
\end{table*}

\begin{table*}[ht]
\centering
\adjustbox{max width=\textwidth}{
\centering
\setlength\tabcolsep{4pt} % control the horizontal spacing between columns
\begin{tabular}{c@{\hspace{5pt}}c@{\hspace{5pt}}c@{\hspace{5pt}}c}
\toprule
Method & Latent NLL - Pinwheel & Latent NLL - Swiss Roll & Latent NLL - Square\\ 
\midrule
AEVB & $1.68 \pm 0.31$ & $5.65 \pm 1.58$ & $2.78 \pm 0.61$ \\
AEVB-IAF & $1.64 \pm 0.73$ & $4.43 \pm 1.09$ & $1.68 \pm 0.27$ \\
AAEB & $-$ & $-$ & $-$ \\
H-AEVB & $2.25 \pm 3.02$ & $4.04 \pm 4.62$ & $2.22 \pm 2.03$ \\
DDVI & $\mathbf{1.27 \pm 0.21}$ & $\mathbf{3.86 \pm 1.17}$ & $\mathbf{1.56 \pm 0.09}$ \\
\textit{DDVI (w/o sleep term)} & $2.12$ & $5.25$ & $2.97$ \\
\bottomrule
\end{tabular}}
\caption{Unsupervised learning on MNIST, including the results of DDVI without the sleep term.}
\vspace{-10pt}
\label{tab:unsupervised_w_failed}
\end{table*}

\section{ELBO for Auxiliary-Variable Generative Models}

We aim to derive the lower bound for the log-likelihood $\log p_\theta(\bx)$ by introducing auxiliary variables and applying the Evidence Lower Bound (ELBO) twice.

Introduce the latent variable $\bz$ and apply the ELBO:
\begin{align}
    \log p_\theta(\bx) 
    &= \log \int p_\theta(\bx, \bz) \, d\bz \\
    &\geq \mathbb{E}_{q_\phi(\bz|\bx)} \left[ \log \frac{p_\theta(\bx, \bz)}{q_\phi(\bz|\bx)} \right] \quad \text{(by Jensen's inequality)} \\
    &= \mathbb{E}_{q_\phi(\bz|\bx)} \left[ \log p_\theta(\bx|\bz) \right] - \text{KL}\left( q_\phi(\bz|\bx) \,\|\, p(\bz) \right) \\
    % &= \mathbb{E}_{q_\phi(\bz|\bx)} \left[ \log p_\theta(\bx|\bz) \right] - \text{KL}\left( q_\phi(\bz|\bx) \,\|\, p_\theta(\bz|\bx) \right)
\end{align}

Then, introduce an auxiliary variable $\by$ and apply the ELBO again:
\begin{align}
    \log p_\theta(\bx) 
    &\geq \mathbb{E}_{q_\phi(\bz|\bx)} \left[ \log p_\theta(\bx|\bz) \right] - \text{KL}\left( q_\phi(\bz|\bx) \,\|\, p(\bz) \right) \\
    &= \mathbb{E}_{q_\phi(\bz|\bx)} \left[ \log p_\theta(\bx|\bz) \right] - \mathbb{E}_{q_\phi(\bz|\bx)} \log q_\phi(\bz|\bx) + \mathbb{E}_{q_\phi(\bz|\bx)} \log p(\bz) \\
    &\geq \mathbb{E}_{q_\phi(\bz|\bx)} \left[ \mathbb{E}_{q_\phi(\by|\bx,\bz)} \left[ \log p_\theta(\bx|\bz) \right] \right] - \mathbb{E}_{q_\phi(\bz|\bx)} \left[ \mathbb{E}_{q_\phi(\by|\bx, \bz)} \left[ \log \frac{q_\phi(\by, \bz|\bx)}{r(\by|\bx, \bz)} \right] \right] + \mathbb{E}_{q_\phi(\bz|\bx)} \left[ \mathbb{E}_{q_\phi(\by|\bx,\bz)} \left[ \log p(\bz) \right] \right] \\
    &= \mathbb{E}_{q_\phi(\by,\bz|\bx)} \left[ \log p_\theta(\bx|\bz) \right] - \text{KL}\left( q_\phi(\by,\bz|\bx) \,\|\, r(\by|\bx,\bz) p(\bz) \right) 
\end{align}

\section{Diffusion Regularization}

We begin with the definition of \(\Lc_\text{diff}(\phi)\), derived as a lower bound on \(\Lc_\text{sleep}(\phi)\):
\[
\Lc_\text{diff}(\phi) = \Exp_{p_\theta(\bx,\bz)} \left[ \Exp_r \left[ \log \frac{q_\phi(\by, \bz | \bx)}{r(\by | \bz, \bx)} \right] \right] + \bar{H}(p_\theta).
\]

To simplify \(\Lc_\text{diff}\), we leverage the Markov structure of the forward process \(r(\by | \bz, \bx)\) and the reverse process \(q_\phi(\by, \bz | \bx)\).

The forward process \(r(\by | \bz, \bx)\) is decomposed as:
\[
r(\by_{1:T} | \bz, \bx) = \prod_{t=1}^T r(\by_t | \by_{t-1}, \bx),
\]
where \(\by_0 = \bz\).

The reverse process \(q_\phi(\by, \bz | \bx)\) is decomposed as:
\[
q_\phi(\by, \bz | \bx) = q_\phi(\by_{0:T} | \bx) = q_\phi(\by_T | \bx) \prod_{t=1}^T q_\phi(\by_{t-1} | \by_t, \bx),
\]
where \(\by_0 = \bz\) and \(\by_{1:T}\) are increasingly noisy versions of \(\bz\).

Substituting the factorizations of \(r\) and \(q_\phi\) into the definition of \(\Lc_\text{diff}(\phi)\), and rewriting the logarithm and rearranging terms, we have:
\begin{align*}
\Lc_\text{diff}(\phi) &= \Exp_{p_\theta(\bx, \bz)} \Bigg[ \Exp_{r(\by_{1:T} | \bz, \bx)} \Bigg[ \log \frac{q_\phi(\by_T | \bx) \prod_{t=1}^T q_\phi(\by_{t-1} | \by_t, \bx)}{\prod_{t=1}^T r(\by_t | \by_{t-1}, \bx)} \Bigg] \Bigg] + \bar{H}(p_\theta) \\
& = \Exp_{r(\by_{1:T}, \bz, \bx)} \left[\log q_\phi(\bz | \by_1, \bx) - \sum_{t=2}^T \KL(r_t || q_t)\right] - \KL(r(\by_T | \bz, \bx) || q_\phi(\by_T | \bx)) + \bar{H}(p_\theta)
\end{align*}

Here, \(\log q_\phi(\bz | \by_1, \bx)\) corresponds to the reconstruction term for the initial latent state \(\bz\), while the summation represents the KL divergence between the forward and reverse processes at each intermediate step. \(r_t = r(\by_{t-1} | \by_t, \bz, \bx)\) is the conditional forward distribution at step \(t\), \(q_t = q_\phi(\by_{t-1} | \by_t, \bx)\) is the reverse process distribution at step \(t\). $\bar H(p_\theta)$ is the expected conditional entropy of $p_\theta(\bz|\bx)$, a constant that does not depend on $\phi$.

This form mirrors the ELBO derivation for diffusion models \citep{sohl2015deep}, where each step in the Markov chain contributes a KL divergence term, and the reconstruction term arises from the connection between the noisy latent variables and the original data.


% Here, we have introduced $q_\phi(\by|\bx,\bz)$ as the variational distribution for the auxiliary variable $\by$.

% Next, we combine the KL divergences involving $\bz$ and $\by$:
% \begin{align}
%     \log p_\theta(\bx) 
%     &\geq \mathbb{E}_{q_\phi(\bz|\bx)} \left[ \mathbb{E}_{q_\phi(\by|\bx,\bz)} \left[ \log p_\theta(\bx|\bz) \right] \right] \nonumber \\
%     &\quad - \text{KL}\left( q_\phi(\bz|\bx) \,\|\, p_\theta(\bz|\bx) \right) \nonumber \\
%     &\quad - \mathbb{E}_{q_\phi(\bz|\bx)} \left[ \text{KL}\left( q_\phi(\by|\bx,\bz) \,\|\, r(\by|\bx,\bz) \right) \right] \label{eqn:elbo_combined}
% \end{align}

% This step leverages the fact that the KL divergence can be decomposed when introducing additional variational distributions.

% Rearrange the terms to express the ELBO in terms of joint distributions:
% \begin{align}
%     \log p_\theta(\bx) 
%     &\geq \mathbb{E}_{q_\phi(\by, \bz|\bx)} \left[ \log p_\theta(\bx|\bz) \right] \nonumber \\
%     &\quad - \text{KL}\left( q_\phi(\by, \bz|\bx) \,\|\, r(\by|\bx,\bz) p(\bz) \right) \label{eqn:elbo_final}
% \end{align}

% Here, we have defined the joint variational distribution $q_\phi(\by, \bz|\bx) = q_\phi(\bz|\bx) q_\phi(\by|\bx,\bz)$ and the joint auxiliary distribution $r(\by|\bx,\bz) p(\bz)$.

% Combining all the steps, we obtain the final ELBO:
% \begin{align}
%     \log p_\theta(\bx) 
%     &\geq \log p_\theta(\bx) - \text{KL}\left( q_\phi(\bz|\bx) \,\|\, p_\theta(\bz|\bx) \right) \label{eqn:elbo2_z1} \\
%     &\geq \log p_\theta(\bx) - \text{KL}\left( q_\phi(\bz|\bx) \,\|\, p_\theta(\bz|\bx) \right) \label{eqn:elbo2_y1}  \\
%     &\quad - \mathbb{E}_{q_\phi(\bz|\bx)} \left[ \text{KL}\left( q_\phi(\by|\bx,\bz) \,\|\, r(\by|\bx,\bz) \right) \right] \nonumber \\
%     &= \mathbb{E}_{q_\phi(\by,\bz|\bx)} \left[ \log p_\theta(\bx|\bz) \right] \label{eqn:elbo2_y2} \\
%     &\quad - \text{KL}\left( q_\phi(\by,\bz|\bx) \,\|\, r(\by|\bx,\bz) p(\bz) \right) \nonumber
% \end{align}

% The derivation shows that by introducing an auxiliary variable $\by$ and applying the ELBO twice, we obtain a tighter lower bound on the log-likelihood $\log p_\theta(\bx)$. The final ELBO incorporates expectations over both $\bz$ and $\by$, as well as KL divergences that measure the divergence between the variational distributions and the corresponding conditional or prior distributions.



% \section{Reproducibility Checklist}
% This paper:
% \begin{itemize}
%     \item Includes a conceptual outline and/or pseudocode description of AI methods introduced: yes
%     \item Clearly delineates statements that are opinions, hypothesis, and speculation from objective facts and results: yes
%     \item Provides well marked pedagogical references for less-familiare readers to gain background necessary to replicate the paper: yes
% \end{itemize}

% \textbf{Does this paper make theoretical contributions? Yes.}

% If yes, please complete the list below.

% \begin{itemize}
%     \item All assumptions and restrictions are stated clearly and formally: yes
%     \item All assumptions and restrictions are stated clearly and formally: yes
%     \item All novel claims are stated formally (e.g., in theorem statements): yes
%     \item Proofs of all novel claims are included: yes
%     \item Proof sketches or intuitions are given for complex and/or novel results: yes
%     \item Appropriate citations to theoretical tools used are given: yes
%     \item All theoretical claims are demonstrated empirically to hold: yes
%     \item All experimental code used to eliminate or disprove claims is included: yes
% \end{itemize}

% \textbf{Does this paper rely on one or more datasets? Yes.}

% If yes, please complete the list below.

% \begin{itemize}
%     \item A motivation is given for why the experiments are conducted on the selected datasets: yes
%     \item All novel datasets introduced in this paper are included in a data appendix: yes
%     \item All novel datasets introduced in this paper will be made publicly available upon publication of the paper with a license that allows free usage for research purposes: yes
%     \item All datasets drawn from the existing literature (potentially including authors' own previously published work) are accompanied by appropriate citations: yes
%     \item All datasets drawn from the existing literature (potentially including authors' own previously published work) are publicly available: yes
%     \item All datasets that are not publicly available are described in detail, with explanation why publicly available alternatives are not scientifically satisfying: yes
% \end{itemize}

% \vspace{12pt}
% \textbf{Does this paper include computational experiments? Yes.}

% If yes, please complete the list below.

% \begin{itemize}
%     \item Any code required for pre-processing data is included in the appendix: yes
%     \item All source code required for conducting and analyzing the experiments is included in a code appendix: yes
%     \item All source code required for conducting and analyzing the experiments will be made publicly available upon publication of the paper with a license that allows free usage for research purposes: yes
%     \item All source code implementing new methods have comments detailing the implementation, with references to the paper where each step comes from: yes
%     \item If an algorithm depends on randomness, then the method used for setting seeds is described in a way sufficient to allow replication of results: yes
%     \item This paper specifies the computing infrastructure used for running experiments (hardware and software), including GPU/CPU models; amount of memory; operating system; names and versions of relevant software libraries and frameworks: yes
%     \item This paper formally describes evaluation metrics used and explains the motivation for choosing these metrics: yes
%     \item This paper states the number of algorithm runs used to compute each reported result: yes
%     \item Analysis of experiments goes beyond single-dimensional summaries of performance (e.g., average; median) to include measures of variation, confidence, or other distributional information: yes
%     \item The significance of any improvement or decrease in performance is judged using appropriate statistical tests (e.g., Wilcoxon signed-rank): yes
%     \item This paper lists all final (hyper-)parameters used for each model/algorithm in the paper's experiments: yes
%     \item This paper states the number and range of values tried per (hyper-)parameter during development of the paper, along with the criterion used for selecting the final parameter setting: yes
% \end{itemize}


\end{document}
